\documentclass{article}
\usepackage{graphicx, amsfonts, amsmath, amssymb, amsthm, float, cancel, braket}
\setcounter{secnumdepth}{3} % Numera solo fino alle subsection (livello 2)
\setcounter{tocdepth}{3}    % Mostra nell'indice fino alle subsubsection (livello 3)

\title{Meccanica Statistica}
\author{Alessandro Ballini}
\date{March 2026}

\begin{document}

\maketitle
\tableofcontents
\newpage

\section*{Introduzione}
In particolare ci occuperemo di sistemi termodinamici all'equilibrio (termico, meccanico e chimico), dunque l'hamiltoniana che descrive i sistemi in esame non dipenderà esplicitamente dal tempo (altrimenti anche il valor medio di una qualsiasi osservabile ne dipenderà, in contrasto con l'ipotesi di equilibrio termodinamico). \\
Lo scopo è quello di ricavare leggi generali a partire dallo studio microscopico di sistemi con un numero molto elevato di gradi di libertà. \\
Inizialmente seguiremo un approccio classico per la descrizione dei fenomeni microscopici. Successivamente tratteremo anche gas quantistici di fermioni e bosoni.


\newpage
\section{Ensemble microcanonico}
\subsection{Densità di probabilità}
Consideriamo un sistema isolato costituito da $N$ particelle, ognuna con $3$ gradi di libertà spaziali. Indichiamo con $q^i_j$ e $p^i_j$ rispettivamente la $j$-esima coordinata spaziale e il $j$-esimo impulso della $i$-esima particella. \\
Fissato un tempo $t_0$, se sono note tutte le coordinate e tutti gli impulsi di ogni particella, lo stato del sistema è identificato univocamente da un punto $(\mathbf{p}(t_0), \mathbf{q}(t_0))$ dello spazio delle fasi ($6N$ dimensionale), dove
\begin{gather}
    \mathbf{q}(t_0) = (q^1_1(t_0), q^1_2(t_0), q^1_3(t_0), q^2_1(t_0), \dots, q^N_3(t_0)) \notag \\
    \mathbf{p}(t_0) = (p^1_2(t_0), p^1_2(t_0), p^1_3(t_0), p^2_1(t_0), \dots, p^N_3(t_0)) \notag
\end{gather}
Inoltre, nota l'hamiltoniana $\mathcal{H}(\mathbf{p}, \mathbf{q})$ che governa il sistema, è possibile evolvere lo stato iniziale nel tempo, dunque si ottiene una curva nello spazio delle fasi, detta \textit{traiettoria di fase}, che indica l'evoluzione nel tempo del sistema, ovvero di tutte le coordinate e tutti gli impulsi. In particolare, dalle equazioni del moto di Hamilton, otteniamo
\begin{gather}
    \begin{cases}
        \dot{q}^i_j &= \frac{\partial \mathcal{H}}{\partial p^i_j} \\
        \dot{p}^i_j &= -\frac{\partial \mathcal{H}}{\partial q^i_j}
    \end{cases} \tag{1.1.1}
\end{gather}
Consideriamo un'osservabile $\theta$ dipendente da $(\mathbf{p}, \mathbf{q})$. Dato che i tempi caratteristici per i fenomeni microscopici sono estremamente più brevi dei tempi caratteristici per i fenomeni globali (o macroscopici), è possibile definire il valore medio di $\theta$ come segue
\begin{gather}
    \braket{\theta} = \lim_{T \rightarrow \infty}\frac{1}{T}\int_0^T \theta(\mathbf{p}(t), \mathbf{q}(t)) \, dt \tag{1.1.2}
\end{gather}
calcolato lungo la traiettoria delle fasi (per semplicità si è posto $t_0 = 0$). Per ipotesi ergotica si ha
\begin{gather}
    \braket{\theta} = \int\theta(\mathbf{p}, \mathbf{q})\rho(\mathbf{p}, \mathbf{q})d\mathbf{p}d\mathbf{q} \tag{1.1.3}
\end{gather}
dove $\rho$ è detta \textit{distribuzione statistica} e rappresenta la densità di probabilità che il sistema si trovi nello stato identificato dal punto $(\mathbf{p}, \mathbf{q})$ dello spazio delle fasi. Ovviamente $\rho$ non può dipendere esplicitamente dal tempo, altrimenti anche il valore medio delle osservabili dipenderebbe dal tempo e ciò non è compatibile con l'ipotesi di equilibrio termodinamico. \\ \\
L'ipotesi ergotica si basa sulla seguente osservazione: consideriamo un elemento di "volume" dello spazio delle fasi $\Delta\mathbf{p}\Delta\mathbf{q}$, dato il comportamento assai complesso della traiettoria delle fasi, è ragionevole supporre che, durante un intervallo di tempo $T$ sufficientemente lungo, la traiettoria delle fasi passi molte volte per ogni elemento di "volume" di questo genere. Sia $\Delta t$ l'intervallo di tempo in cui la traiettoria delle fasi si trova all'interno di $\Delta\mathbf{p}\Delta\mathbf{q}$, allora si ha che la quantità
\begin{gather}
    \omega := \lim_{T \rightarrow \infty} \frac{\Delta t}{T} \tag{1.1.4}
\end{gather}
è finita e può essere interpretata come la probabilità di trovare il sistema in un stato contenuto nel "volume" $\Delta\mathbf{p}\Delta\mathbf{q}$. \\
Se il "volume" considerato diventa infinitesimo, ovvero si considera
\begin{gather}
    d\mathbf{p}d\mathbf{q} = dp^1_1dq^1_1 \cdot \dots \cdot dp^N_3dq^N_3 \notag
\end{gather}
allora si può introdurre la probabilità infinitesima $d\omega$ data da
\begin{gather}
    d\omega = \rho(\mathbf{p}, \mathbf{q}) \, d\mathbf{p}d\mathbf{q} \tag{1.1.5}
\end{gather}
dove $\rho$ è una PDF (definita per ogni coppia $(\mathbf{p}, \mathbf{q})$ nello spazio delle fasi), ovviamente normalizzata. \\ \\
Vediamo come $\rho$ evolve nel tempo. Consideriamo una regione $\Omega$ dello spazio delle fasi e sia $\Sigma = \partial \Omega$. La probabilità $P$ di trovare il sistema in un stato contenuto in $\Omega$ è data da
\begin{gather}
    P = \int_{\Omega}\rho \, d\mathbf{p}d\mathbf{q} \tag{1.1.6}
\end{gather}
Il sistema è isolato, quindi non ci sono sorgenti o pozzi. Posto $\mathbf{v} = (\dot{\mathbf{p}}, \dot{\mathbf{q}})$, segue che $P$ soddisfa la relazione di continuità
\begin{gather}
    \frac{\partial P}{\partial t} = -\int_{\Sigma} \rho\mathbf{v} \cdot \hat{\boldsymbol{n}} \, d\mathbf{p}d\mathbf{q} \tag{1.1.7}
\end{gather}
con $\hat{\boldsymbol{n}}$ versore ortogonale (uscente) alla (iper)superficie $\Sigma$. La relazione globale $(1.1.7)$ è equivalente alla seguente relazione locale
\begin{align}
    \frac{\partial\rho}{\partial t} = -\nabla \cdot (\rho\mathbf{v}) \tag{1.1.8}
\end{align}
Infatti, supponendo $\rho$ è sufficientemente regolare, si ha
\begin{align}
    \frac{\partial P}{\partial t} &= \frac{\partial}{\partial t}\int_\Omega \rho \, d\mathbf{p}d\mathbf{q} \notag \\
    &= \int_\Omega \frac{\partial\rho}{\partial t} \, d\mathbf{p}d\mathbf{q} \notag
\end{align}
quindi si ottiene
\begin{gather}
    \int_\Omega \frac{\partial \rho}{\partial t} \, d\mathbf{p}d\mathbf{q} = -\int_\Omega \rho\mathbf{v} \cdot \hat{\boldsymbol{n}} \, d\mathbf{p}d\mathbf{q} \tag{1.1.9}
\end{gather}
ovvero
\begin{gather}
    \int_\Sigma \left(\frac{\partial \rho}{\partial t} + \rho\mathbf{v} \cdot \hat{\boldsymbol{n}}\right) \, d\mathbf{p}d\mathbf{q} = 0 \tag{1.1.10}
\end{gather}
dall'arbitrarietà di $\Omega$, segue che la funzione integranda in $(1.1.10)$ è identicamente nulla (ovunque per continuità, ovvero anche su insiemi di misura nulla), pertanto concludiamo
\begin{gather}
    \frac{\partial \rho}{\partial t} + \nabla \cdot (\rho\mathbf{v}) = 0 \tag{1.1.11}
\end{gather}
Esplicitiamo la divergenza di $\rho\mathbf{v}$ ricordando che $\mathbf{v} = (\dot{\mathbf{p}}, \dot{\mathbf{q}})$ soddisfa la equazioni di Hamilton lungo la traiettoria delle fasi. In particolare, sia $N$ il numero di particelle, si ha
\begin{align}
    \frac{\partial \rho}{\partial t} + \nabla\cdot(\rho\mathbf{v}) &= \frac{\partial \rho}{\partial t} + \sum_{i = 1}^{3N}\rho\left(\frac{\partial\dot{p}_i}{\partial p_i} + \frac{\partial\dot{q}_i}{\partial q_i}\right) + \left(\dot{p}_i\frac{\partial\rho}{\partial p_i} + \dot{q}_i\frac{\partial\rho}{\partial q_i}\right) \notag \\
    &= \frac{\partial \rho}{\partial t} + \sum_{i = 1}^{3N}\rho\cancel{\left(-\frac{\partial^2\mathcal{H}}{\partial p_i \partial q_i} + \frac{\partial^2\mathcal{H}}{\partial q_i\partial p_i}\right)} + \left(-\frac{\partial\mathcal{H}}{\partial q_i}\frac{\partial\rho}{\partial p_i} + \frac{\partial \mathcal{H}}{\partial p_i}\frac{\partial\rho}{\partial q_i}\right) \notag \\
    &= \frac{\partial \rho}{\partial t} + \{\rho, \mathcal{H}\} = 0 \tag{1.1.12}
\end{align}
Come detto in precedenza, se il sistema si trova all'equilibrio termodinamico deve per forza valere $\frac{\partial\rho}{\partial t} = 0$, pertanto $\rho$ soddisfa $\{\rho, \mathcal{H}\} = 0$. \\
Tuttavia, $(1.1.12)$ è ancora di un'equazione differenziale impossibile da risolvere per $N$ molto grande. Si semplifica molto se è possibile scrivere $\rho$ come funzione di $\mathcal{H}(\mathbf{p}, \mathbf{q})$ stesso, ovvero se è possibile scrivere $\rho = \rho(\mathcal{H}(\mathbf{p}, \mathbf{q}))$. In tal caso
\begin{align}
    \{\rho, \mathcal{H}\} &= \frac{\partial \rho}{\partial \mathcal{H}}\sum_{i = 1}^{3N}\cancel{\left(-\frac{\partial \mathcal{H}}{\partial q_i}\frac{\partial \mathcal{H}}{\partial p_i} + \frac{\partial \mathcal{H}}{\partial p_i}\frac{\partial \mathcal{H}}{\partial q_i}\right)} \tag{1.1.13}
\end{align}
ovvero ogni termine della somma è identicamente nullo. \\ \\
Abbiamo essenzialmente dimostrato che l'equazione di continuità per $\rho$ equivale a $\frac{d\rho}{dt} = 0$. Segue che la distribuzione di probabilità $\rho$ è uniforme lungo tutta la (iper)superficie definita dall'energia (fissata) del sistema. Poniamo
\begin{gather}
    \Gamma(E) = \int_{\Sigma(E)} d\mathbf{p}d\mathbf{q} \tag{1.1.14}
\end{gather}
la misura della (iper)superficie $\Sigma(E)$ definita da $\mathcal{H}(\mathbf{p}, \mathbf{q}) = E$. $\Gamma(E)$ è proporzionale al "numero" di stati $(\mathbf{p}, \mathbf{q})$ che soddisfano $\mathcal{H}(\mathbf{p}, \mathbf{q}) = E$, con fattore di proporzionalità dato dalla misura minima dell'elemento infinitesimo nello spazio delle fasi, ovvero $1/h^{3N}$ ($\Delta p_i\Delta q_i \geq h$ per il \textit{principio di indeterminazione}). \\
$\rho$ deve essere normalizzata su $\Sigma(E)$, pertanto
\begin{gather}
    \rho = \frac{1}{\Gamma(E)} \tag{1.1.15}
\end{gather}
Per comodità si introduce un parametro $\Delta > 0$ e si decide arbitrariamente che $\rho$ è costante lungo il volume identificato dall'equazione $E \leq H(\mathbf{p}, \mathbf{q}) \leq E + \Delta$, in modo tale da non dover trattare $\rho$ come una distribuzione. Tale assunzione è valida affinché l'introduzione di questo parametro non comporta effetti critici nel sistema.  Allora si ha
\begin{gather}
    \Gamma(E) = \int_{E \leq H(\mathbf{p}, \mathbf{q}) \leq E + \Delta} d\mathbf{p}d\mathbf{q} \tag{1.1.16}
\end{gather}
Ogni sistema isolato descritto dalle variabili termodinamiche $E, N, V$ appertiene ad un ensemble detto \textit{ensemble microcanonico}. \\

\newpage
\subsection{Entropia}
Consideriamo un ensemble microcanonico con energia fissata $E$, definiamo l'entropia del sistema come segue
\begin{gather}
    S = k\log \Gamma(E) \tag{1.2.1}
\end{gather}
con $k$ costante di Boltzmann. Vogliamo dimostrare che questa definizione di entropia è estensiva e quindi può essere usata come potenziale termodinamico. Inoltre, vogliamo dimostrare che è legata alla definizione macroscopica di entropia data da
\begin{gather}
    TdS = dE + PdV - \mu dN \tag{1.2.2}
\end{gather}
Consideriamo un sistema isolato confinato in un volume $V = V_1 + V_2$ e con un numero $N = N_1 + N_2$ di particelle indipendenti. \\
Pensiamo di dividere il sistema in due sottosistemi $1$ e $2$ a volume e numero di particelle fissato ($V_1, V_2$ e $N_1, N_2$), pertanto l'unica interazione sarà termica ($dV = 0 = dN$). Fissiamo l'energia totale del sistema $E$ e siano $E_1$ ed $E_2$ rispettivamente l'energia del sottosistema $1$ e l'energia del sottosistema $2$. $E_1$ ed $E_2$ non sono quantità fissate, sono libere di variare affinché rispettino il vincolo $E = E_1 + E_2$. Siano $S_1 = k\log \Gamma_1(E_1)$ e $S_2 = k\log \Gamma_2(E_2)$ rispettivamente l'entropia sottosistema $1$ e l'entropia del sottosistema $2$. Allora
\begin{gather}
    \Gamma(E) = \sum_{i = 1}^{E/\Delta} \Gamma_1(E^{(i)}_1)\Gamma_2(E - E^{(i)}_1) \tag{1.2.3}
\end{gather}
dove $E/\Delta$ è il numero di combinazioni possibili, $\Delta$ rappresenta la spaziatura tra i livelli energetici (caso non degenere). Dato che $\Gamma(E)$ è essenzialmente proporzionale al numero di stati, $\Gamma_1(E_1)$ e $\Gamma_2(E_2)$ si combinano col prodotto, ovvero prendiamo ogni combinazione possibile di stati (che rispetta il vincolo). Poiché che si tratta di un numero finito di termini, esiste $\overline{E}_1$ tale che
\begin{gather}
    \Gamma_1(\overline{E}_1)\Gamma_2(E - \overline{E}_1) \geq \Gamma_1(E^{(i)}_1)\Gamma_2(E - E^{(i)}_1) \tag{1.2.4}
\end{gather}
per ogni $i = 1, \dots, E/\Delta$. Possiamo minorare e maggiorare $\Gamma(E)$ come segue
\begin{gather}
    \Gamma_1(\overline{E}_1)\Gamma_2(E - \overline{E}_1) \leq \Gamma(E) \leq \left(\frac{E}{\Delta}\right) \Gamma_1(\overline{E}_1)\Gamma_2(E - \overline{E}_1) \tag{1.2.5}
\end{gather}
Dalla \textit{teoria cinetica dei gas} sappiamo che $E = \frac{3}{2}NkT$, dunque $E/\Delta \sim N$, mentre la misura della (iper)superficie degli stati ad energia $E$, che in questo caso è una sfera nello spazio degli impulsi, va come $R^{3N - 1}$, con $R = \sqrt{2mE}$. Segue che 
\begin{align*}
    \Gamma(E) \sim E^{\frac{3N - 1}{2}} 
\end{align*}
Concludiamo che il termine $E/\Delta$ è trascurabile per $N \rightarrow \infty$ (limite termodinamico) e da $(1.2.5)$ si ottiene
\begin{gather}
    \Gamma_1(\overline{E}_1)\Gamma_2(E - \overline{E}_1) \leq \Gamma(E) \leq \Gamma_1(\overline{E}_1)\Gamma_2(E - \overline{E}_1) \notag \qquad N \rightarrow \infty
\end{gather}
ovvero
\begin{gather}
    \Gamma(E) = \Gamma_1(\overline{E}_1)\Gamma_2(E - \overline{E}_1) \qquad N \rightarrow \infty \tag{1.2.6}
\end{gather}
Pertanto, dato che la funzione $\log$ è monotona e quindi preserva la disuguaglianza, al limite termodinamico si ha
\begin{gather}
    k\log\Gamma(E) = k\log\Gamma_1(\overline{E}_1) + k\log\Gamma_2(E - \overline{E}_1) \tag{1.2.7}
\end{gather}
che equivale a
\begin{gather}
    S(E) = S_1(\overline{E}_1) + S_2(E - \overline{E}_1) \tag{1.2.8}
\end{gather}
Per determinare $\overline{E}_1$, e quindi anche $\overline{E}_2 = E - \overline{E}_1$, cerchiamo per quale $E_1$ vale
\begin{gather}
    \frac{\partial}{\partial E_1}\Gamma_1(E_1)\Gamma_2(E_2) = 0 \tag{1.2.9}
\end{gather}
ovvero
\begin{align}
    \Gamma_2(E_2)\frac{\partial}{\partial E}\Gamma_1(E_1) + \Gamma_1(E_1)\frac{\partial E_2}{\partial E_1}\frac{\partial}{\partial E_2}\Gamma_2(E_2) = 0 \tag{1.2.10}
\end{align}
In particolare
\begin{gather}
    \frac{\partial E_2}{\partial E_1} = \frac{\partial}{\partial E_1}(E - E_1) = -1 \notag
\end{gather}
quindi la condizione che deve essere soddisfatta è la seguente
\begin{gather}
    \Gamma_2\frac{\partial\Gamma_1}{\partial E_1} = \Gamma_1\frac{\partial\Gamma_2}{\partial E_2} \tag{1.2.11}
\end{gather}
Dividendo entrambi i membri per $\Gamma_1\Gamma_2$ e esprimendo in termini di logaritmi, $(1.2.11)$ diventa
\begin{gather}
    \frac{\partial}{\partial E_1}\log \Gamma_1 = \frac{\partial}{\partial E_2}\log \Gamma_2 \tag{1.2.12}
\end{gather}
A questo punto ricordiamo che per $dV = 0 = dN$ si ha $TdS = dE$, che all'equilibrio termico, cioè $T_1 = T_2$, comporta
\begin{gather}
    \frac{\partial S_1}{\partial E_1} = \frac{\partial S_2}{\partial E_2} \tag{1.2.13}
\end{gather}
confrontando $(1.2.13)$ e $(1.2.12)$ viene naturale pensare che $S$ e $\log\Gamma$ sono legate da qualche costante di proporzionalità $k$, come avevamo già ipotizzato, cioè
\begin{align*}
    S = k\log\Gamma
\end{align*}
Vediamo come sono legate la misura della "corteccia" definita dall'energia del sistema e la misura della sfera (piena) contenuta la suo interno. Poniamo
\begin{align}
    \Omega(E) &= \int_{\mathcal{H} \leq E}d\mathbf{p}d\mathbf{q} \tag{1.2.14}
\end{align}
dalla definizione $(1.1.16)$ segue
\begin{gather}
    \Gamma(E) = \Omega(E + \Delta) - \Omega(E) \notag
\end{gather}
Per $\Delta$ sufficientemente piccolo possiamo scrivere
\begin{gather}
    \Gamma(E) \sim \Delta\frac{\partial \Omega}{\partial E} \tag{1.2.15}
\end{gather}
In particolare $\Omega \sim E^\alpha$, con $\alpha \sim N$, dato che si tratta del "volume" della (iper)sfera. Sostituendo la $(1.2.15)$ nella $(1.2.1)$ si ottiene
\begin{align}
    k\log\Gamma &\sim k\log\left(\Delta\frac{\partial\Omega}{\partial E}\right) \notag \\
    &= k\log\left(\frac{\Delta}{E}\right) + k\log\left(E\frac{\partial\Omega}{\partial E}\right) \notag \\
    &= k\log\left(\frac{\Delta}{E}\right) + k\log\left(\alpha\Omega\right) \notag \\
    &\sim k\log\Omega \notag
\end{align}
dove è stato trascurato il primo termine del membro di destra per $N \rightarrow \infty$, concludiamo che, al limite termodinamico, possiamo sostituire $\Gamma$ con $\Omega$.

\newpage
\subsection{Teorema di equipartizione dell'energia}
\textbf{Teorema (di equipartizione dell'energia per ensemble microcanonico):} sia $\mathcal{H}$ un'hamiltoniana associata ad un sistema microcanonico, sia $x_i$, con $i = 1, \dots, 6N$, uno qualsiasi tra $(p_1, \dots, p_{3N}, q_1, \dots, q_{3N})$, allora si ha 
\begin{align}
    \braket{x_i\frac{\partial \mathcal{H}}{\partial x_j}} = \delta_{ij}kT \tag{1.3.1}
\end{align}
per ogni $i, j = 1, \dots, 6N$. \\

\textbf{Dimostrazione:} per un ensemble microcanico la media di una qualsiasi osservabile $A$ è data da 
\begin{align}
    \braket{A} = \int A(\mathbf{p}, \mathbf{q})\rho(\mathbf{p}, \mathbf{q}) \, d\mathbf{p}d\mathbf{q} \tag{1.3.2}
\end{align}
dove $\rho$ è la funzione densità di probabilità. In particolare $\rho$ è costante lungo la "corteccia" identificata da $E \leq \mathcal{H} \leq E + \Delta$, ovvero 
\begin{align}
    \rho(\mathbf{p}, \mathbf{q}) = \begin{cases}
        \frac{1}{\Gamma(E)} \qquad &E \leq \mathcal{H}(\mathbf{p}, \mathbf{q}) \leq E + \Delta \\
        0 \qquad &\text{altrimenti}
    \end{cases} \tag{1.3.3}
\end{align}
Allora 
\begin{align}
    \braket{A} = \frac{1}{\Gamma(E)}\int_{E \leq \mathcal{H} \leq E + \Delta}A(\mathbf{p}, \mathbf{q}) \, d\mathbf{p}d\mathbf{q} \tag{1.3.4}
\end{align}
dunque, sostituendo in $(1.3.3)$, si ottiene
\begin{align}
    \braket{x_i\frac{\partial \mathcal{H}}{\partial x_j}} &= \frac{1}{\Gamma(E)}\int_{E \leq \mathcal{H} \leq E + \Delta} x_i\frac{\partial \mathcal{H}}{\partial x_j} \, d\mathbf{p}, d\mathbf{q} \notag \\
    &= \frac{1}{\Gamma(E)}\left(\int_{\mathcal{H} \leq E + \Delta}x_i\frac{\partial \mathcal{H}}{\partial x_j} \, d\mathbf{p}d\mathbf{q} - \int_{E \leq \mathcal{H}}x_i\frac{\partial \mathcal{H}}{\partial x_j} \, d\mathbf{p}d\mathbf{q}\right) \notag \\
    &\sim \frac{\Delta}{\Gamma(E)}\frac{\partial}{\partial E}\int_{E \leq \mathcal{H}}x_i\frac{\partial \mathcal{H}}{\partial x_j} \, d\mathbf{p}d\mathbf{q} \notag \\
    &= \left(\frac{\partial \Omega}{\partial E}\right)^{-1}\frac{\partial}{\partial E}\int_{E \leq \mathcal{H}}x_i\frac{\partial \mathcal{H}}{\partial x_j} \, d\mathbf{p}d\mathbf{q} \notag \\
    &= \left(\frac{\partial \Omega}{\partial E}\right)^{-1}\frac{\partial}{\partial E}\int_{E \leq \mathcal{H}}x_i\frac{\partial}{\partial x_j}(\mathcal{H} - E) \, d\mathbf{p}d\mathbf{q} \notag \\
    &= \left(\frac{\partial \Omega}{\partial E}\right)^{-1}\frac{\partial}{\partial E}\left(\cancel{\int_{E \leq \mathcal{H}}\frac{\partial}{\partial x_j}(x_i(\mathcal{H} - E)) \, d\mathbf{p}d\mathbf{q}} - \delta_{ij}\int_{E \leq \mathcal{H}}(\mathcal{H} - E) \, d\mathbf{p}d\mathbf{q}\right) \notag
\end{align}
Il primo integrale è nullo dato che $x_i(\mathcal{H} - E) = 0$ per $\mathcal{H} = E$, ovvero lungo il bordo dell'insieme di integrazione. Allora si ha 
\begin{align}
    \braket{x_i\frac{\partial \mathcal{H}}{\partial x_j}} &= \delta_{ij}\left(\frac{\partial \Omega}{\partial E}\right)^{-1}\frac{\partial}{\partial E}\int_{E \leq \mathcal{H}}(E - \mathcal{H}) \, d\mathbf{p}d\mathbf{q} \notag \\
    &= \delta_{ij}\left(\frac{\partial\Omega}{\partial E}\right)^{-1}\int_{E \leq \mathcal{H}}\frac{\partial}{\partial E}(E - \mathcal{H}) \, d\mathbf{p}d\mathbf{q} + \cancel{(E - \mathcal{H})|_{\mathcal{H} = E}} \notag \\
    &= \delta_{ij}\Omega\left(\frac{\partial \Omega}{\partial E}\right)^{-1} \notag \\
    &= \delta_{ij}\left(\frac{\partial \log\Omega}{\partial E}\right)^{-1} \notag \\
    &= \delta_{ij}k\left(\frac{\partial S}{\partial E}\right)^{-1} \notag \\
    &= \delta_{ij}kT \notag
\end{align}
dove si è usata la formula di derivazione di funzione integrale rispetto ad un parametro e $T = \left(\frac{\partial S}{\partial E}\right)^{-1}$. \qed

\newpage
\subsection{Particelle libere}
Vediamo ora il caso particolare del gas libero confinato in un volume $V$. L'hamiltoniana che descrive il sistema è
\begin{gather}
    \mathcal{H} = \sum_{i = 1}^N \frac{\mathbf{p}^2_i}{2m} \notag
\end{gather}
Calcoliamo $\Omega(E)$, si ha
\begin{gather}
    \Omega(E) = \int_{\mathcal{H} \leq E} d\mathbf{p}_1d\mathbf{p}_2 \dots d\mathbf{p}_Nd\mathbf{q}_1d\mathbf{q}_2 \dots d\mathbf{q}_N \tag{1.4.1}
\end{gather}
$\Omega$ è proporzionale al numero di stati contenuti nell'insieme identificato da $\mathcal{H} \leq E$. Dividendo per la misura minima dell'elemento $d\mathbf{p}_1d\mathbf{q}_i$ otteniamo il numero di stati, dunque, a meno di ridefinire $\Omega$, possiamo scrivere
\begin{gather}
    \Omega(E) = \frac{1}{h^{3N}}\int_{\mathcal{H} \leq E}d\mathbf{p}_1d\mathbf{p}_2 \dots d\mathbf{p}_Nd\mathbf{q}_1d\mathbf{q}_2 \dots d\mathbf{q}_N \tag{1.4.2}
\end{gather}
L'hamiltoniana $\mathcal{H}$ non dipende da $\mathbf{q}$, dunque possiamo subito calcolare il contributo dell'integrale rispetto alle coordinate $\mathbf{q}$ in $(1.4.2)$. Ogni $d\mathbf{q}_i$ contribuisce con un fattore $V$, ovvero il volume del sistema fisico in esame, pertanto si ha
\begin{gather}
    \Omega(E) = \left(\frac{V}{h^3}\right)^N\omega_{3N}(R) \tag{1.4.3}
\end{gather}
dove $\omega_{3N}(R)$ è la misura della sfera $3N$-dimensionale di raggio $R$. Dobbiamo quindi determinare la seguente quantità per $n$ generico
\begin{gather}
    \omega_n(R) = \int_{\|\mathbf{x}\| \leq R}dx_1 \dots dx_n \notag
\end{gather}
Ragionando per induzione, sappiamo che la misura della sfera bidimensionale di raggio $R$ è $4\pi R^2$ e la misura della sfera tridimensionale di raggio $R$ è $\frac{4}{3}\pi R^3$. Supponiamo che anche la misura della sfera $n - 1$ dimensionale sia proporzionale a $n - 1$, ovvero che esista $c_{n - 1} > 0$ tale che $m(\mathbf{S}^{n - 1}) = c_{n - 1}R^{n - 1}$ e verifichiamo se vale anche per la sfera $n$ dimensionale $\mathbf{S}^n$. $\mathbf{S}^n$ è identificata dalla seguente equazione
\begin{gather}
    \mathbf{S}^n = \{\mathbf{x} \in \mathbf{R}^n : x^2_1 + \dots + x^2_n \leq R^2\} \notag
\end{gather}
Allora
\begin{align}
    \omega_n(R) &= \int_{-R}^{R}\int_{x^2_1 + \dots + x^2_{n - 1} \leq R^2 - x^2_n}dx_1 \dots dx_{n - 1}dx_n \notag \\
    &= \int_{-R}^R c_{n - 1}(R^2 - x^2_n)^{\frac{n - 1}{2}} \, dx_n \notag \\
    &= c_{n - 1}R^n\int_{-1}^1(1 - u^2)^{\frac{n - 1}{2}} \, du \tag{1.4.4}
\end{align}
dove si è posto $u = x_n/R$. A questo punto, possiamo determinare $c_{n - 1}$ notando che, per ipotesi induttiva, si ha
\begin{align*}
    dx_1 \dots dx_{n - 1} = (n - 1)c_{n - 1}r^{n - 2}dr
\end{align*}
dove $r = \|\mathbf{x}\|$. Inoltre vale la seguente identità
\begin{align}
    \pi^{\frac{n - 1}{2}} &= \int_{-\infty}^\infty dx_1 \dots \int_{-\infty}^\infty dx_{n - 1}e^{-(x^2_1 + \dots + x^2_{n - 1})} \notag \\
    &= (n - 1)c_{n - 1}\int_{0}^\infty r^{n - 2}e^{-r^2} \, dr \notag \\
    &= \left(\frac{n}{2} - \frac{1}{2}\right)c_{n - 1}\int_0^\infty t^{\frac{n}{2} - \frac{3}{2}}e^{-t} \, dt \notag \\
    &= \left(\frac{n}{2} - \frac{1}{2}\right)c_{n - 1}\Gamma\left(\frac{n}{2} - \frac{1}{2}\right) \notag \\
    &= c_{n - 1}\Gamma\left(\frac{n}{2} + \frac{1}{2}\right) \notag
\end{align}
ovvero
\begin{gather}
    c_{n - 1} = \frac{\pi^{\frac{n - 1}{2}}}{\Gamma\left(\frac{n + 1}{2}\right)} \tag{1.4.5}
\end{gather}
Pertanto si ottiene
\begin{align}
    \omega_n(R) &= \frac{\pi^{\frac{n - 1}{2}}}{\Gamma\left(\frac{n + 1}{2}\right)}R^n\int_{-1}^1(1 - u^2)^{\frac{n - 1}{2}} \, du \notag \\
    &= \frac{\pi^{\frac{n}{2}}}{\Gamma\left(\frac{n}{2} + 1\right)}R^n \tag{1.4.6}
\end{align}
Sostituendo in $(1.4.3)$ si ha
\begin{gather}
    \Omega(E) = c_{3N}\left(\frac{V}{h^3}(2mE)^{3/2}\right)^N \tag{1.4.6}
\end{gather}
A questo punto sostituendo nella definizione di entropia otteniamo
\begin{gather}
    S(E, N, V) = k\left(\log c_{3N} + N\log \frac{V}{h^3} + \frac{3}{2}\log (2mE)\right) \tag{1.4.7}
\end{gather}
per poter trattare l'espressione $(1.4.7)$ dell'entropia $S$ come potenziale termodinamico dobbiamo far uso della \textit{formula di Stirling}, ovvero
\begin{gather}
    \log\Gamma\left(\frac{3N}{2} + 1\right) \sim \frac{3N}{2}\log\frac{3N}{2} - \frac{3N}{2} \qquad N \rightarrow \infty \notag
\end{gather}
allora
\begin{gather}
    \log c_{3N} \sim \frac{3N}{2}\log\pi - \frac{3N}{2}\log\frac{3N}{2} + \frac{3N}{2} \qquad N \rightarrow \infty \notag
\end{gather}
e
\begin{gather}
    S(E, N, V) \sim kN\log\left(V\left(\frac{4\pi m}{3h^2}\frac{E}{N}\right)^{3/2}\right) + \frac{3}{2}Nk \qquad N \rightarrow \infty \tag{1.4.8} 
\end{gather}
Esprimendo $E$ in funzione di $(N, S, V)$ e usando la relazione macroscopica $TdS = dE + PdV - \mu dN$, possiamo ricavare le relazione che lega i potenziali termodinamici con le grandezze termodinamiche globali, infatti
\begin{gather}
    E(N, S, V) = \left(\frac{3h^2}{4\pi m}\right)\frac{N}{V^{2/3}}\exp\left(\frac{2}{3}\frac{S}{Nk} - 1\right) \tag{1.4.9}
\end{gather}
e in particolare
\begin{align}
    T &= \left(\frac{\partial E}{\partial S}\right)_{N, V} \notag \\
    &= \frac{2}{3Nk}E \notag
\end{align}
ovvero
\begin{gather}
    E = \frac{3}{2}NkT \notag
\end{gather}
pertanto ritroviamo la relazione dell'energia interna del gas ricavata con la teoria cinetica dei gas e comprendiamo che il termine $3/2$ deriva dal fatto che l'hamiltoniana è quadratica negli impulsi e che il sistema è tridimensionale. Inoltre
\begin{align}
    P &= -\frac{\partial E}{\partial V} \notag \\
    &= \frac{2}{3}\frac{E}{V} \notag \\
    &= \frac{NkT}{V} \notag
\end{align}
che comporta
\begin{gather}
    PV = NkT \notag
\end{gather}
ovvero ricaviamo la legge dei gas ideali. \\ \\
Tuttavia, l'espressione $(1.4.8)$ dell'entropia non è estensiva a causa del termine $V$ all'interno del logaritmo. Questo accade perché non abbiamo tenuto conto dell'indistinguibilità classica delle particelle per il gas ideale nell'integrazione rispetto a $d\mathbf{q}$ per il calcolo di $\Omega(E)$. Necessitiamo dunque di un fattore $N!$, detto \textit{fattore di Gibbs}, a dividere $\Gamma$ (o $\Omega$) per correggere il problema. Tale fattore apparirà autonomamente quando considereremo l'indistinguibilità quantistica delle particelle. \\ \\
Poniamo $u = E/N$ e $v = V/N$ e riscriviamo $S$ come segue
\begin{align}
    S(E, N, V) = kN\log(Vu^{2/3}) + \frac{3}{2}kNS_0 \tag{1.4.10}
\end{align}
dove
\begin{align}
    S_0 = 1 + \log\frac{4\pi m}{3h^2} \notag
\end{align}
Il fattore $V$ all'interno del logaritmo comporta un problema. Affinché l'entropia $S$ sia correttamente estensiva deve essere possibile fattorizzare la dipendenza da $N$ in modo tale che $S$ possa essere scritta il prodotto tra $N$ e una funzione $s$ che dipende solamente da variabili intensive, ovvero
\begin{align}
    S = Ns(E/N, V/N) \notag
\end{align}
Aggiungendo a mano il fattore di Gibbs $N!$ si ha
\begin{align}
    k\log\Omega \longrightarrow k\log\Omega - k\log N! \tag{1.4.11}
\end{align}
e al limite termodinamico, usando la formula di Stirling, si ottiene
\begin{align}
    k\log N! &= k\log\Gamma(N + 1) \notag \\
    &\sim kN\log N - N \notag
\end{align}
Pertanto, l'espressione corretta dell'entropia diventa
\begin{align}
    S(E, N, V) &= kN\log\left(\frac{V}{N}\left(\frac{E}{N}\right)^{3/2}\right) + \frac{5}{2}kNS_0 \notag \\
    &= kN\log(uv^{3/2}) - \frac{5}{2}kNS_0 \tag{1.4.12}
\end{align}
notiamo che $N$ è già fattorizzato e i termini che moltiplica dipendono solamente da quantità intensive, dunque questa espressione dell'entropia è estensiva. 

\newpage
\subsection{Particelle soggette a potenziale armonico}
Nel caso in cui le particelle sono soggette ad un potenziale armonico si ha
\begin{align}
    \mathcal{H} = \sum_{i = 1}^N \left(\frac{\mathbf{p}^2_i}{2m} + \frac{1}{2}m\omega^2\mathbf{q}^2_i\right) \tag{1.5.1}
\end{align}
Al limite termodinamico possiamo approssimare
\begin{align}
    k\log\Gamma \sim k\log\Omega \notag
\end{align}
Allora
\begin{align}
    \Omega = \frac{1}{h^{3N}} \int_{\mathcal{H} \leq E} d\mathbf{p}d\mathbf{q} \notag
\end{align}
Poniamo
\begin{align}
    \begin{cases}
        \mathbf{x}_i = \frac{\mathbf{p}_i}{2m} \qquad &i = 1, \dots, N \\
        \mathbf{y}_i = \frac{1}{2}m\omega^2\mathbf{q}_i \qquad &i = 1, \dots, N
    \end{cases} \notag
\end{align}
Allora
\begin{align}
    \Omega &= \frac{1}{h^{3N}}\left(2m\right)^{\frac{3N}{2}}\left(\frac{2}{m\omega^2}\right)^{\frac{3N}{2}}\int_{\sum \|\mathbf{x}_i\|^2 + \|\mathbf{y}_i\|^2 \leq E} d\mathbf{x}_1 \dots d\mathbf{x}_Nd\mathbf{y}_1 \dots d\mathbf{y}_N \notag \\
    &= \frac{1}{h^{3N}}\left(2m\right)^{\frac{3N}{2}}\left(\frac{2}{m\omega^2}\right)^{\frac{3N}{2}}\frac{\pi^{3N}}{3N\Gamma(3N)}E^{3N} \notag \\
    &= \left(\frac{2}{h\omega}\right)^{3N}\frac{\pi^{3N}}{3N\Gamma(3N)}E^{3N} \notag \\
    &= \left(\frac{E}{\hbar\omega}\right)^{3N}\frac{1}{(3N)!} \notag
\end{align}
quindi, usando la formula di Stirling al limite termodinamico, si ha
\begin{align}
    S &\sim k\log \Omega \notag \\
    &\sim -3kN\log 3N + 3kN + 3kN\log \frac{E}{\hbar\omega} \notag \\
    &= 3kN\left(1 + \log\left(\frac{1}{3\hbar\omega}\frac{E}{N}\right)\right) \tag{1.5.2}
\end{align}
Notiamo che l'entropia è correttamente estensiva, infatti $N$ è fattorizzato e il termine che moltiplica dipende solamente da quantità intensive, in questo caso $u = E/N$ densità di energia. Questo segue dal fatto che le particelle in questo caso sono distinguibili, infatti sono ben identificate dal punto di equilibrio rispetto al quale oscillano. Inoltre
\begin{align}
    \frac{1}{T} &= \left(\frac{\partial S}{\partial E}\right)_{V, N} \notag \\
    &= \frac{3kN}{E} \notag
\end{align}
ovvero
\begin{align}
    E = 3NkT \tag{1.5.3} \notag
\end{align}
Il fattore $3$ dipende dal fatto che in questo caso i gradi di libertà sono $6$, $3$ per gli impulsi e $3$ per le coordinate, e l'hamiltoniana è quadratica in entrambi le variabili dinamiche. Come vedremo successivamente questo fenomeno è una conseguenza dell'equipartizione dell'energia sui gradi di libertà. \\
Inoltre notiamo che
\begin{align}
    P &= T\left(\frac{\partial S}{\partial V}\right)_{E, N} = 0 \notag
\end{align}
infatti, essendo ogni particella localizzata attorno ad un punto di equilibrio, non avvengono urti che si manifestano poi in pressione.

\newpage
\subsection{Particelle ultra-relativistiche libere}
L'hamiltoniana di un insieme di $N$ particelle libere ultra-relativistiche in una scatola è
\begin{align}
    \mathcal{H} = c\sum_{i = 1}^N |\mathbf{p}_i| \tag{1.6.1}
\end{align}
Al limite termodinamico si ha $\Gamma \sim \Omega$, dunque siamo interessati a calcolare il seguente integrale
\begin{align}
    \Omega = \frac{1}{h^{3N}}\int_{\sum |\mathbf{p}_i| \leq E/c} d\mathbf{p}d\mathbf{q} \tag{1.6.2}
\end{align}
L'equazione
\begin{align}
    \sum_{i = 1}^N|\mathbf{p}_i| \leq \frac{E}{c} \tag{1.6.3}
\end{align}
identifica un rombo nello spazio $N$ dimensionale dei moduli degli impulsi. In questo caso, dato che $|\mathbf{p}_i| = \sqrt{p^2_{ix} + p^2_{iy} + p^2_{iz}}$, è conveniente passare in coordinate polari, allora si ha
\begin{align}
    \Omega = \left(\frac{4\pi V}{h^3}\right)^N \int_{\sum p_i \leq E/c} \prod_{i = 1}^N p^2_i \, dp_i \tag{1.6.4}
\end{align}
dove $p_i = |\mathbf{p}_i|$, pertanto $p_i \geq 0$ per ogni $i = 1, \dots, N$. Determiniamo in generale
\begin{align}
    I_n(A) = \int_{\sum x_i \leq A}\prod_{i = 1}^nx^2_i \, dx_i \tag{1.6.5}
\end{align}
dove $x_i \geq 0$ per ogni $i = 1, \dots, n$. Per $n = 2$ si ha
\begin{align}
    I_2(A) &= \int_{x_1 + x_2 \leq A}x^2_1x^2_2 \, dx_1dx_2 \notag \\
    &= \int_0^A x^2_1 \int_0^{A - x_1} x^2_2 \, dx_2dx_1 \notag \\
    &= \frac{1}{3}\int_0^A x^2_1(A - x_1)^3 \, dx_1 \tag{1.6.6}
\end{align}
Poniamo
\begin{align}
    J_k(A) = \int_0^A t^2(A - t)^k \, dt \notag
\end{align}
si ha che
\begin{align}
    J_k(A) &= \int_0^A t^2(A - t)^k \, dt \notag \\
    &= \int_0^A (A - s)^2s^k \, ds \notag \\
    &= \int_0^A (A^2s^k + s^{k + 2} - 2As^{k + 1}) \, ds \notag \\
    &= A^{k + 3}\left(\frac{1}{k + 1} + \frac{1}{k + 3} - \frac{2}{k + 2}\right) \notag \\
    &= A^{k + 3}\frac{2k!}{(k + 3)!} \tag{1.6.7}
\end{align}
Pertanto, $(1.6.6)$ diventa
\begin{align}
    I_2(A) = \frac{1}{3}J_3(A) = \frac{2^2}{6!}A^6 \tag{1.6.8}
\end{align}
Per $n = 3$ si ha
\begin{align}
    I_3(A) = \int_{x_1 + x_2 + x_3 \leq A}x^2_1x^2_2x^2_3 \, dx_1dx_2dx_3 &= \int_0^A x^2_1 \int_{x_2 + x_3 \leq A - x_1}x^2_2x^2_3 \, dx_2dx_3dx_1 \notag \\
    &= \int_0^A x^2_1I_2(A - x_1) \, dx_1 \notag \\
    &= \frac{2^2}{6!}\int_0^A x^2_1(A - x_1)^6 \notag \\
    &= \frac{2^2}{6!}J_6(A) \notag \\
    &= \frac{2^3}{9!}A^9 \tag{1.6.9}
\end{align}
Confrontando $(1.6.8)$ e $(1.6.9)$ viene naturale ragionare per induzione, dunque supponiamo che per $k < n$ valga
\begin{align}
    I_k(A) = \frac{2^k}{(3k)!}A^{3k} \notag
\end{align}
e dimostriamo che vale anche per $k + 1$. Si ha che
\begin{align}
    I_{k + 1}(A) &= \int_{x_1 + \dots +x_{k + 1} \leq A}x^2_1 \dots x^2_{k + 1} \, dx_1 \dots dx_{k + 1} \notag \\
    &= \int_0^A x^2_1 \int_{x_2 + \dots + x_{k + 1} \leq A - x_1}x^2_2 \dots x^2_{k + 1} \, dx_1 \dots dx_{k + 1}dx_1 \notag \\
    &= \int_0^A x^2_1I_k(A - x_1) \, dx_1 \notag \\
    &= \frac{2^k}{(3k)!}\int_0^A x^2_1(A - x_1)^{3k} \, dx_1 \notag \\
    &= \frac{2^k}{(3k)!}J_{3k}(A) \notag \\
    &= \frac{2^k}{(3k)!}\frac{2(3k)!}{(3(k + 1))!}A^{3(k + 1)} \notag \\
    &= \frac{2^{k + 1}}{(3(k + 1))!}A^{3(k + 1)} \notag
\end{align}
che è la tesi, dunque abbiamo dimostrato che
\begin{align}
    I_n(A) = \frac{2^n}{(3n)!}A^{3n} \tag{1.6.10}
\end{align}
per ogni $n \in \mathbf{N}$. A questo punto, $(1.6.4)$ diventa
\begin{align}
    \Omega = \left(\frac{4\pi V}{h^3}\right)^N\frac{2^N}{(3N)!}\left(\frac{E}{c}\right)^{3N} \tag{1.6.11}
\end{align}
In questo caso dobbiamo aggiungere a mano il fattore di Gibbs per far sì che l'entropia sia correttamente estensiva, dunque
\begin{align}
    \Omega = \frac{1}{N!}\frac{1}{(3N)!}\left(\frac{8\pi V E^3}{h^3c^3}\right)^N \tag{1.6.12}
\end{align}
Allora
\begin{align}
    S &\sim k\log \Omega \notag \\
    &\sim kN\left(4 + \log\left(\frac{8\pi VE^3}{27h^3c^3N^4}\right)\right) \tag{1.6.13}
\end{align}
dove si è usata nuovamente la formula di Stirling al limite termodinamico. Notiamo che
\begin{align}
    \frac{1}{T} &= \left(\frac{\partial S}{\partial E}\right)_{V, N} \notag \\
    &= \frac{3kN}{E} \notag
\end{align}
ovvero
\begin{align}
    E = 3NkT \tag{1.6.14}
\end{align}
in questo caso il fattore $\frac{1}{2}$ non è presente perché l'hamiltoniana non è quadratica.


\newpage
\subsection{Sistema di $N$ siti a due stati}
Consideriamo un insieme di $N$ siti localizzati a due stati, in particolare ogni sito del sistema può avere energia $0$ o energia $\varepsilon$ fissata. Fissata l'energia del sistema siamo interessati a determinare l'espressione dell'entropia al limite termodinamico. \\ \\
Fissiamo l'energia totale del sistema $E = m\varepsilon$, dove $m$ è il numero di stati ad energia $\varepsilon$. Il primo sito con energia $\varepsilon$ ha $N$ possibili opzioni. Fissato il primo, il secondo ha $N - 1$ possibili opzioni e così via fino ad aver assegnato ogni quanto di energia $\varepsilon$. Tuttavia l'ordine di scelta non cambia il sistema finale, finché gli stati con quanto di energia non nullo siano sempre gli stessi. Pertanto
\begin{align}
    \Gamma(m, N) = \frac{N!}{(N - m)!m!} \tag{1.7.1}
\end{align}
Al limite termidinamico si ha 
\begin{align}
    S(m, N) &\sim k\left(N\log N - N - (N - m)\log(N - m) + (N - m) - m\log m + m\right) \notag \\
    &= k\left(N\log N - (N - m)\log(N - m) - m\log m\right) \tag{1.7.2}
\end{align}
Notiamo che 
\begin{align}
    \frac{1}{T} &= \left(\frac{\partial S}{\partial E}\right)_N \notag \\
    &= \left(\frac{\partial m}{\partial E}\frac{\partial S}{\partial m}\right)_N \notag \\
    &= \frac{k}{\varepsilon}\log 
    \left(\frac{N}{m} - 1\right) \tag{1.7.3}
\end{align}
ovvero 
\begin{align}
    e^{\frac{\varepsilon}{kT}} &= \frac{N}{m} - 1 \tag{1.7.4}
\end{align}
Possiamo quindi esprimere $m$ in funzione della temperatura 
\begin{align}
    m = \frac{N}{e^{\frac{\varepsilon}{kT}} + 1} \tag{1.7.5}
\end{align}
Notiamo che per $T \rightarrow \infty$ si ha $m = N/2$, ovvero quando l'energia termica è molto più grande del quanto di energia $\varepsilon$ si osserva la configurazione per la quale $\Gamma$ è massimo. Mentre per $T \rightarrow 0$ si ha $m = 0$, ovvero a basse temperature non ci sono quanti di energia accessibili.

\newpage
\subsection{Sistema di $N$ oscillatori armonici quantistici} 
Consideriamo un sistema composto da $N$ oscillatori armonici quantistici $3D$. Fissata l'energia del sistema, siamo interessati a determinare l'espressione dell'entropia al limite termodinamico. \\ \\
Limitiamoci per un attimo al caso di oscillatori armonici quantistici unidimensionali. Possiamo pensare a tale sistema come un insieme di $N$ scatole localizzate, ognuna con un numero $n_i$ di quanti di energia. In particolare, fissata l'energia totale del sistema $E$, si ha 
\begin{align}
    E = \sum_{i = 1}^N(n_i + \frac{1}{2}) = \sum_{i = 1}^N n_i + \frac{N}{2} \tag{1.8.1}
\end{align}
dove si è espressa l'energia di ogni oscillatore armonico quantistico in unità di $\hbar\omega$. Il sistema è costruito da una successione (finita) di oggetti, ovvero i quanti di energia e le pareti delle scatole. A parte le pareti "esterne" che sono fissate, ci sono $N - 1$ pareti e $M = E - \frac{N}{2}$ quanti disponibili, abbiamo quindi un insieme di $M + N - 1$ oggetti. Tuttavia, non sono tutti indistinguibili, infatti i quanti formano un sottoinsieme composto da oggetti tutti indindistinguibili tra loro e lo stesso vale per le pareti. Concludiamo che il numero totale di configurazioni distinguibili è 
\begin{align}
    \Gamma(N, M) &= \frac{(M + N - 1)!}{M!(N - 1)!} \tag{1.8.2}
\end{align}
Questa formula vale in generale, infatti dato un insieme $U$ con cardinalità $N$ che può essere espresso come unione di $k$ sottoinsiemi $V_1, \dots, V_k$ composti da oggetti indistinguibili, ognuno con cardinalità $n_i$, $i = 1, \dots, k$, si ha che il numero $N_\sigma$ di permutazioni $\sigma$ è dato da
\begin{align}
    N_\sigma = \frac{N!}{n_1!n_2!\dots n_k!} \notag
\end{align}
detto coefficiente multinomiale. \\
Nel caso $3D$ basta sostituire $N$ con $3N$, infatti l'hamiltoniano dell'oscillatore armonico quantistico $3D$ è la somma di 3 hamiltoniani unidimensionali. Allora si ha 
\begin{align}
    \Gamma(N, M) &= \frac{(M + 3N - 1)!}{M!(3N - 1)!} \tag{1.8.3}
\end{align}
Al limite termodinamico $3N - 1 \sim 3N$, quindi, usando la formula di Stirling, si ottiene 
\begin{align}
    S(N, M) &\sim k\left((M + 3N)\log(M + 3N) - M\log M - 3N\log 3N\right) \notag
\end{align}
esplicitando la dipendenza dall'energia totale $E$ e fattorizzando $N$ si ottiene 
\begin{align}
    S(E, N) &\sim kN\left(\left(u + \frac{3}{2}\right)\log N\left(u + \frac{3}{2}\right) - \left(u - \frac{3}{2}\right)\log N\left(u - \frac{3}{2}\right) - 3\log 3N\right) \notag \\
    &= kN\left(u\log \frac{u + 3/2}{u - 3/2} + \frac{3}{2}\log\left(u^2 - \frac{9}{4}\right) -3\log3\right) \tag{1.8.4}
\end{align}
Notiamo che 
\begin{align}
    \frac{1}{T} &= \left(\frac{\partial S}{\partial E}\right)_N \notag \\
    &= \left(\frac{\partial u}{\partial E}\frac{\partial S}{\partial u}\right)_N \notag \\
    &= k\log \left(\frac{u + 3/2}{u - 3/2}\right) \notag
\end{align}
ovvero 
\begin{align}
    e^{\frac{1}{kT}} &= \frac{u + 3/2}{u - 3/2} \tag{1.8.5}
\end{align}
quindi si ottiene 
\begin{align}
    u(T) &= \frac{3}{2}\frac{e^{\frac{1}{kT}} + 1}{e^{\frac{1}{kT}} - 1} \notag \\
    &= \frac{3}{2}\frac{e^{\frac{1}{2kT}} + e^{-\frac{1}{2kT}}}{e^{\frac{1}{2kT}} - e^{-\frac{1}{2kT}}} \notag \\
    &= \frac{3}{2}\coth\frac{1}{2kT} \tag{1.8.6}
\end{align}
dove si è moltiplicato e diviso per $e^{-\frac{1}{2kT}}$ il secondo membro della prima uguaglianza. Moltiplicando per la scala scelta all'inizio (unità di $\hbar\omega$) si ha 
\begin{align}
    u(T) = \frac{3}{2}\hbar\omega\coth\frac{\hbar\omega}{2kT} \tag{1.8.7}
\end{align}
Notiamo che $u \rightarrow \infty$ per $T \rightarrow \infty$, in particolare 
\begin{align}
    u(T) \sim \frac{3}{2}kT \tag{1.8.8}
\end{align}
per $T \rightarrow \infty$, ovvero, al limite in cui l'energia termica è molto maggiore di $\hbar\omega$, si osserva comportamento classico e l'informazione quantistica viene persa. Mentre per $T \rightarrow 0$ si ha 
\begin{align}
    u(T) \sim \frac{3}{2}\hbar\omega \tag{1.8.9}
\end{align}
ovvero $E = \frac{3}{2}N\hbar\omega$ (in unità di $\hbar\omega$), quindi ogni oscillatore armonico si trova allo stato fondamentale $n = 0$.

\newpage
\section{Ensemble canonico}
Siamo interessati a descivere statisticamente un sistema a contatto con un reservoir. Possiamo immaginare un reservoir come un sistema ad energia e volume molto maggiori rispetto al sistema in esame, che quindi resta imperturbato nonostante l'interazione esterna.

\subsection{Energia libera}
Nel contesto dell'ensemble microcanonico, consideriamo un sistema $X$ descritto dalla terna $(E, N, V)$ e consideriamo un sottoinsieme $X_1$ di $X$ descritto dalla terna $(E_1, N_1, V_1)$. Poniamo $X_2 = X \setminus X_1$, il complementare di $X_1$, descritto dalla terna $(E_2, N_2, V_2)$ e supponiamo che $X_2$ sia molto più grande di $X_1$. Sotto tali ipotesi, l'interazione tra $X_1$ e $X_2$, che è essenzialmente limitata ad una regione limitrofa al bordo di $X_1$, diventa infinitesima, e quindi trascurabile, al limite termidinamico, ovvero $V_1$ molto grande e $V >> V_1$. Possiamo dunque trattare $X_1$ come un sistema isolato. \\
Fissate l'energia totale $E$ e l'energia $E_1$ di $X_1$, l'energia $E_2$ di $X_2$ deve soddisfare 
\begin{align}
    E_2 = E - E_1 \notag
\end{align}
La densità di probabilità $\rho$ associata a $X$ è costante sull'ipersuperficie identificata da $\mathcal{H}(\mathbf{p}, \mathbf{q}) = E$ e nulla altrove. In particolare 
\begin{align}
    \rho(\mathbf{p}, \mathbf{q}) = \begin{cases}
        \frac{1}{\Gamma(E)} \qquad &\mathcal{H}(\mathbf{p}, \mathbf{q}) = E \\
        0 \qquad &\text{altrimenti}
    \end{cases} \notag
\end{align}
La probabilità di misurare energia $E_1$ in $X_1$ è proporzionale al numero di microstati di $X_2$ compatibili con il vincolo $E = E_1 + E_2$, ovvero 
\begin{align}
    \rho_1(\mathbf{p}^{(1)}, \mathbf{q}^{(1)}) \propto \Gamma_2(E - E_1) \tag{2.1.1}
\end{align}
Poiché $E_1 << E$ possiamo sviluppare in serie di Taylor 
\begin{align}
    \Gamma_2(E - E_1) \sim \Gamma_2(E) - E_1\frac{\partial \Gamma_2}{\partial E_2} \tag{2.1.2}
\end{align}
e dalla definizione statistica dell'entropia del sistema si ottiene  
\begin{align}
    \rho_1(\mathbf{p}^{(1)}, \mathbf{q}^{(1)}) &\propto \exp\left(\frac{S_2(E)}{k}-\frac{E_1}{kT}\right) \notag
\end{align}
dove si è usato $\frac{\partial S_2}{\partial E_2} = \frac{1}{T_2} \sim \frac{1}{T}$. In particolare, il termine $S_2(E)$ non dipende dalle variabili microscopiche del sottosistema $X_1$, pertanto si ottiene 
\begin{align}
    \rho_1(\mathbf{p}^{(1)}, \mathbf{q}^{(1)}) &\propto \exp\left(-\frac{\mathcal{H}_{N_1}(\mathbf{p}^{(1)}, \mathbf{q}^{(1)})}{kT}\right) \tag{2.1.3}
\end{align}
dove si è esplicitata la dipendenza dall'hamiltoniana dal numero di particelle. \\ \\
In generale, dato un sistema descritto da $(E, N, V)$ posto a contatto con un reservoir si ha 
\begin{align}
    \rho(\mathbf{p}, \mathbf{q}) \propto \exp\left(-\frac{\mathcal{H}_N(\mathbf{p}, \mathbf{q})}{kT}\right) \tag{2.1.4}
\end{align}
Resta da determinare la costante di normalizzazione, che poi, come nel paragrafo precedente, determinerà il potenziale termodinamico associato alla descrizione del sistema. Definiamo 
\begin{align}
    Q_N = \frac{1}{h^{3N}}\int \exp\left(-\frac{\mathcal{H}_N(\mathbf{p}, \mathbf{q})}{kT}\right) \, d\mathbf{p}d\mathbf{q} \tag{2.1.5}
\end{align}
Allora 
\begin{align}
    \frac{1}{h^{3N}Q_N}\int \rho(\mathbf{p}, \mathbf{q}) \, d\mathbf{p}d\mathbf{q} = 1 \tag{2.1.6}
\end{align}
Notiamo che $Q_N = Q_N(T, V)$, ovvero dipende dalla terna $(N, T, V)$. Il potenziale termodinamico le cui variabili naturali sono $(N, T, V)$ è l'\textit{energia libera di Helmholtz} $F$, infatti 
\begin{align}
    F = E - TS \notag
\end{align}
pertanto 
\begin{align}
    dF &= dE - SdT - TdS \notag \\
    &= dE - SdT -dE - PdV + \mu dN \notag \\
    &= -SdT - PdV + \mu dN \tag{2.1.7}
\end{align}
Facendo riferimento all'espressione $(2.1.5)$, diamo come definizione statistica dell'energia libera la seguente 
\begin{align}
    F(N, T, V) = -kT\log Q_N(T, V) \tag{2.1.8}
\end{align}
Verifichiamo che sia correttamente estensiva. Poiché per ipotesi $X_1$ e $X_2$ sono debolmente interagenti si ha che $\mathcal{H}_{tot} = \mathcal{H}_1 + \mathcal{H}_2$, allora 
\begin{align}
    Q_N &= \left(\int\frac{d\mathbf{p}^{(1)}d\mathbf{q}^{(1)}}{h^{3N_1}}e^{-\frac{\mathcal{H}_1(\mathbf{p}^{(1)}, \mathbf{q}^{(1)})}{kT}}\right)\left(\int\frac{d\mathbf{p}^{(2)}d\mathbf{q}^{(2)}}{h^{3N_2}}e^{-\frac{\mathcal{H}_2(\mathbf{p}^{(2)}, \mathbf{q}^{(2)})}{kT}}\right) \notag \\
    &= Q^{(1)}_{N_1}Q^{(2)}_{N_2} \notag
\end{align}
segue che 
\begin{align}
    F_{tot} &= -kT\log Q_N \notag \\
    &= -kT\log(Q^{(1)}_{N_1}Q^{(2)}_{N_2}) \notag \\
    &= F_1 + F_2 \tag{2.1.9}
\end{align}
ovvero $F$ correttamente estensiva (sotto l'ipotesi di sottosistemi debolmente interagenti). \\
Osserviamo che il termine $E$ nella definizione macroscopica dell'energia libera non è l'energia del sistema, infatti quest'ultima è una variabile stocastica per gli ensemble canonici, dato che il sistema scambia energia con il reservoir per termalizzarsi. La corretta definizione di $F$ è 
\begin{align}
    F = \braket{\mathcal{H}_N} - TS \tag{2.1.10}
\end{align}
Per mostrarlo poniamo $\beta = \frac{1}{kT}$, si ha che 
\begin{align}
    e^{\beta F}Q_N &= e^{\beta F}\int\frac{d\mathbf{p}d\mathbf{q}}{h^{3N}}e^{-\beta\mathcal{H}_N} = 1 \notag
\end{align}
ovvero 
\begin{align}
    \int\frac{d\mathbf{p}d\mathbf{q}}{h^{3N}}e^{\beta(F - \mathcal{H}_N)} = 1 \tag{2.1.11}
\end{align}
allora 
\begin{align}
    \frac{\partial}{\partial \beta}\int\frac{d\mathbf{p}d\mathbf{q}}{h^{3N}}e^{\beta(F - \mathcal{H})} &= \int\frac{d\mathbf{p}d\mathbf{q}}{h^{3N}}\left(F - \mathcal{H}_N + \beta\left(\frac{\partial F}{\partial \beta}\right)_{N, V}\right)e^{\beta(F - \mathcal{H}_N)} \notag \\
    &= \left(F + \beta\left(\frac{\partial F}{\partial \beta}\right)_{N, V}\right) - \frac{1}{Q_N}\int\frac{d\mathbf{p}d\mathbf{q}}{h^{3N}}\mathcal{H}_Ne^{-\beta\mathcal{H}_N} \notag \\
    &= \left(F - T\left(\frac{\partial F}{\partial T}\right)_{N, V}\right) - \braket{\mathcal{H}_N} = 0 \notag
\end{align}
Inoltre, da $(2.1.7)$ si ha
\begin{align}
    -T\left(\frac{\partial F}{\partial T}\right)_{N, V} = TS \notag 
\end{align}
pertanto concludiamo che la corretta definizione di $F$ per ensemble canonici è la $(2.1.10)$.

\newpage
\subsection{Teorema di equipartizione dell'energia}
\textbf{Teorema (di equipartizione dell'energia per ensemble canonico):} sia $\mathcal{H}$ un'hamiltoniana associata ad un sistema canonico, sia $x_i$, $i = 1, \dots, 6N$, uno qualsiasi tra $(p_1, p_2, \dots, p_{3N}, q_1, q_2, \dots, q_{3N})$, allora si ha 
\begin{align}
    \braket{x_i\frac{\partial\mathcal{H}}{\partial x_j}} = \delta_{ij}kT \tag{2.2.1}
\end{align}
per ogni $i, j = 1, \dots, 6N$. \\

\textbf{Dimostrazione:} per l'ensemble canonico si ha 
\begin{align}
    \braket{x_i\frac{\partial\mathcal{H}}{\partial x_j}} &= \frac{1}{Q_Nh^{3N}}\int d\mathbf{p}d\mathbf{q} \, x_i\frac{\partial\mathcal{H}}{\partial x_j}e^{-\beta\mathcal{H}} \notag \\
    &= \frac{1}{Q_Nh^{3N}\beta}\int d\mathbf{p}d\mathbf{q} \, x_i\frac{\partial}{\partial x_j}e^{-\beta\mathcal{H}} \notag \\
    &= \frac{1}{Q_Nh^{3N}\beta}\left(\cancel{\int \prod_{\tilde{x}_j \in (\mathbf{p}, \mathbf{q}) \setminus x_j}d\tilde{x}_j \, x_i e^{-\beta\mathcal{H}}|_{x_j^{(1)}}^{x_j^{(2)}}} + \int d\mathbf{p}d\mathbf{q} \, \frac{\partial x_i}{\partial x_j}e^{-\beta\mathcal{H}}\right) \notag \\
    &= \delta_{ij}kT \notag
\end{align}
poiché la prima funzione integranda si annulla agli estremi (per hamiltoniane note come particella libera e oscillatore armonico).

\newpage
\subsection{Fluttuazioni dell'energia}
Vediamo ora come si comportano le fluttuazioni dell'energia rispetto alla media. Notiamo che la definizione di momento di ordine due (rispetto a $\rho$) dell'hamiltoniana $\mathcal{H}$ è legata al concetto di varianza della distribuzione dell'energia, infatti si ha 
\begin{align}
    \Delta \mathcal{H} = \sqrt{\braket{\mathcal{H}^2} - \braket{\mathcal{H}}^2} \tag{2.3.1}
\end{align}
Consideriamo la media statistica sull'ensemble dell'hamiltoniana $\mathcal{H}$
\begin{align}
    \braket{\mathcal{H}} &= \frac{1}{h^{3N}Q_N}\int d\mathbf{p}d\mathbf{q} \, \mathcal{H}e^{-\beta\mathcal{H}} \notag \\
    &= \frac{1}{h^{3N}}e^{\beta F}\int d\mathbf{p}d\mathbf{q} \, \mathcal{H}e^{-\beta\mathcal{H}} \notag
\end{align}
L'energia libera $F$ non dipende dalle variabili microscopiche del sistema $(\mathbf{p}, \mathbf{q})$, pertanto 
\begin{align}
    \braket{\mathcal{H}} = \frac{1}{h^{3N}}e^{\beta F}\int d\mathbf{p}d\mathbf{q} \, \mathcal{H}e^{-\beta\mathcal{H}} = \frac{1}{h^{3N}}\int d\mathbf{p}d\mathbf{q} \, \mathcal{H}e^{\beta(F - \mathcal{H})} \tag{2.3.2}
\end{align}
per $(2.1.11)$ si ha 
\begin{align}
    \braket{\mathcal{H}} = 1 \cdot \braket{\mathcal{H}} = \left(\frac{1}{h^{3N}}\int d\mathbf{p}d\mathbf{q} \, e^{\beta(F - \mathcal{H})}\right)\braket{\mathcal{H}} \notag
\end{align}
Inoltre, dato che il valor medio dell'hamiltoniana non dipende dalle variabili microscopiche $(\mathbf{p}, \mathbf{q})$, si ottiene 
\begin{align}
    \braket{\mathcal{H}} = \frac{1}{h^{3N}}\int d\mathbf{p}d\mathbf{q} \, \braket{\mathcal{H}}e^{\beta(F - \mathcal{H})} \tag{2.3.3}
\end{align}
Confrontando $(2.3.2)$ e $(2.3.3)$ si ottiene 
\begin{align}
    \frac{1}{h^{3N}}\int d\mathbf{p}d\mathbf{q} \, (\braket{\mathcal{H}} - \mathcal{H})e^{\beta(F - \mathcal{H})} = 0 \tag{2.3.4}
\end{align}
Ricaviamo un vincolo sulla varianza derivando rispetto a $\beta$, infatti 
\begin{align}
    0 &= \frac{\partial}{\partial \beta}\left(\frac{1}{h^{3N}}\int d\mathbf{p}d\mathbf{q} \, (\braket{\mathcal{H}} - \mathcal{H})e^{\beta(G - \mathcal{H})}\right) \notag \\
    &= \frac{1}{h^{3N}}\int d\mathbf{p}d\mathbf{q} \, \left(\frac{\partial \braket{\mathcal{H}}}{\partial \beta} + (\braket{\mathcal{H}} - \mathcal{H})\left(G - \mathcal{H} + \beta\left(\frac{\partial G}{\partial \beta}\right)_{N, V}\right)\right)e^{\beta(G - \mathcal{H})} \tag{2.3.5}
\end{align}
Il termine $\frac{\partial \braket{\mathcal{H}}}{\partial \beta}$ non dipende dalle variabili microscopiche $(\mathbf{p}, \mathbf{q})$, quindi può essere fattorizzato fuori dall'integrale (che fa 1). Inoltre, da $(2.1.10)$ notiamo che 
\begin{align}
    F + \beta\left(\frac{\partial F}{\partial \beta}\right)_{N, V} &= F - T\left(\frac{\partial F}{\partial T}\right)_{N, V} \notag \\
    &= F + TS \notag \\
    &= \braket{\mathcal{H}} \tag{2.3.6}
\end{align}
e sostituendo in $(2.3.5)$ si ha 
\begin{align}
    \frac{\partial \braket{\mathcal{H}}}{\partial \beta} + \braket{(\braket{\mathcal{H}} - \mathcal{H})^2} = 0 \tag{2.3.7}
\end{align}
ponendo $U = \braket{\mathcal{H}}$ si ha 
\begin{align}
    -kT^2U + \braket{(U - \mathcal{H})^2} = 0 \notag
\end{align}
ovvero 
\begin{align}
    \braket{(U - \mathcal{H})^2} = kT^2U \tag{2.3.8}
\end{align}
Dalla termodinamica ricordiamo che 
\begin{align}
    \frac{\partial U}{\partial T} = c_V \tag{2.3.9}
\end{align}
dove $c_V$ è la funzione di risposta capacità termica (a volume costante), deduciamo quindi che le fluttuazioni dell'energia sono strettamente legata alla capacità termica. \\ \\

\newpage
\subsection{Particelle libere}
Consideriamo un sistema composto da $N$ particelle confinate in un volume $V$, in particolare 
\begin{align}
    \mathcal{H}(\mathbf{p}, \mathbf{q}) = \sum_{i = 1}^N \frac{|\mathbf{p}_i|^2}{2m} \notag
\end{align}
dove $|\mathbf{p}_i|^2 = p^2_{ix} + p^2_{iy} + p^2_{iz}$. Siamo interessati a determinare l'espressione dell'energia libera, pertanto calcoliamo la funzione di partizione canonica $Q_N$ associata a $\mathcal{H}$, si ha 
\begin{align}
    Q_N &= \frac{1}{h^{3N}}\int d\mathbf{p}d\mathbf{q} \, e^{-\beta\mathcal{H}} \notag \\
    &= \frac{1}{h^{3N}}\int d\mathbf{p}d\mathbf{q} \, e^{-\beta\sum_{i = 1}^N \frac{|\mathbf{p}_i|^2}{2m}} \tag{2.4.1}
\end{align}
Notiamo che si tratta dell'integrale del prodotto di $N$ gaussiane indipendenti tra loro, pertanto possiamo integrare ogni termine singolarmente
\begin{align}
    Q_N &= \frac{V^N}{N!}\frac{1}{h^{3N}}\prod_{i = 1}^N \int d\mathbf{p}_i e^{-\beta\frac{|\mathbf{p}_i|^2}{2m}} \tag{2.4.2}
\end{align}
dove è stato aggiunto "a mano" anche il fattore di Gibbs necessario per rendere l'energia libera correttamente estensiva. Otteniamo quindi 
\begin{align}
    Q_N = \frac{V^N}{N!}\left(\frac{2\pi mkT}{h^2}\right)^{3N/2} \tag{2.4.3}
\end{align}
Notiamo che 
\begin{align}
    Q_N = \frac{Q^N_1}{N!} \tag{2.4.4}
\end{align}
questo accade ogni qual volta l'hamiltoniana è separabile, ovvero è somma di hamiltoniane di particelle indipendenti tra loro (e indistinguibili, per questo il fattore di Gibbs). \\ \\
Nota la funzione di partizione canonica, possiamo ricavare l'espressione dell'energia libera e studiare l'aspetto termodinamico del sistema. In particolare 
\begin{align}
    F &= -kT\log Q_N \notag \\
    &= -kT\log \frac{V^N}{N!}\left(\frac{2\pi mkT}{h^2}\right)^{3N/2} \notag
\end{align}
usando la formula di Stirling al limite termodinamico si ottiene 
\begin{align}
    F \sim -NkT\log\frac{V}{N} - NkT\left(1 - \frac{3}{2}\log\frac{2\pi mkT}{h^2}\right) \tag{2.4.5}
\end{align}
$F$ è correttamente estensiva, infatti $N$ è già fattorizzato e $F$ può essere scritta come il prodotto tra $N$ ed una funzione che dipende solamente da quantità intensive (in questo caso la densità di volume e la temperatura). \\ \\
Infine calcoliamo l'energia interna del sistema $U = \braket{\mathcal{H}}$. Per definizione di media statistica sull'ensemble canonico
\begin{align}
    U = \frac{1}{Q_Nh^{3N}}\int d\mathbf{p}d\mathbf{q} \, \mathcal{H}e^{-\beta\mathcal{H}} \tag{2.4.6}
\end{align}
Notiamo che 
\begin{align}
    U = -\frac{\partial}{\partial \beta}\log Q_N = -\frac{1}{Q_N}\frac{\partial Q_N}{\partial\beta} \tag{2.4.7}
\end{align}
portando l'operazione di derivazione sotto il segno di integrale. Pertanto 
\begin{align}
    U &= -\frac{\partial}{\partial\beta}\left(-N\log\frac{V}{N} -N\left(1 - \frac{3}{2}\log\frac{2\pi m}{h^2\beta}\right)\right) \notag \\
    &= -\frac{\partial}{\partial\beta}\left(-\frac{3}{2}N\log\frac{1}{\beta}\right) = \frac{3}{2}NkT \tag{2.4.8}
\end{align}
risultato che è in accordo con il \textit{teorema di equipartizione di energia}.

\newpage
\subsection{Particelle soggette a potenziale armonico}
Consideriamo ora un sistema composto da $N$ oscillatori armonici tridimensionali, ovvero 
\begin{align}
    \mathcal{H}(\mathbf{p}, \mathbf{q}) = \sum_{i = 1}^N \frac{|\mathbf{p_i|^2}}{2m} + \frac{1}{2}m\omega^2|\mathbf{q}_i|^2 \notag
\end{align}
Tale sistema è equivalente a un sistema di $3N$ oscillatori armonici unidimensionali, infatti i moti lungo i tre assi si scompongono e l'hamiltoniana di ogni singola particella è separabile e scrivibile come somma di tre hamiltoniane unidimensionali. \\ \\
Anche in questo caso vale la proprietà $(2.4.4)$ dato che gli stati delle singole particelle sono indipendenti. Tuttavia, non è necessario il fattore di Gibbs, infatti le particelle sono distinguibili grazie al termine potenziale dell'hamiltoniana. Per determinare $Q_N$ determiniamo prima $Q_1$. Si ha 
\begin{align}
    Q_1 &= \frac{1}{h^3}\int d\mathbf{p}_1d\mathbf{q}_1 \, e^{-\beta\mathcal{H}_1} \notag \\
    &= \frac{1}{h^3}\int dp_xdp_ydp_zdxdydx \, e^{-\beta\left(\frac{p^2_x + p^2_y + p^2_z}{2m} + \frac{1}{2}m\omega^2(x^2 + y^2 + z^2)\right)} \notag \\
    &= \frac{1}{h^3}\left(\frac{2\pi m}{\beta}\right)^{3/2}\left(\frac{2\pi}{m\omega^2\beta}\right)^{3/2} \notag \\
    &= \left(\frac{kT}{\hbar\omega}\right)^3 \tag{2.5.1}
\end{align}
e quindi 
\begin{align}
    Q_N = \left(\frac{kT}{\hbar\omega}\right)^{3N} \tag{2.5.2}
\end{align}
Segue che 
\begin{align}
    F = -3NkT\log \frac{kT}{\hbar\omega} \tag{2.5.3}
\end{align}
Inoltre 
\begin{align}
    U &= \frac{\partial F}{\partial \beta} \notag \\
    &= -\frac{\partial}{\partial \beta}\log Q_N \notag \\
    &= 3NkT \tag{2.5.4}
\end{align}
come ci si aspetta dal \textit{teorema di equipartizione dell'energia}.

\newpage
\subsection{Sistema di $N$ siti a due stati}
Consideriamo un sistema composto da $N$ siti localizzati a due stati. In particolare, ogni sito può avere energia $E_+ = +\varepsilon$ o $E_- = -\varepsilon$ (con $\varepsilon > 0)$. \\
Se l'energia è fissata (contesto dell'ensemble microcanonico) sappiamo come è fatta la funzione $\Gamma(E)$ e quindi sappiamo determinare l'espressione dell'entropia microcanonica, ovvero, posto $N_+$ il numero di siti con energia $+\varepsilon$ e $N_-$ il numero di siti con energia $-\varepsilon$, si ha 
\begin{align}
    \Gamma(E) = \frac{N!}{N_+!N_-!} \tag{2.6.1}
\end{align}
dove $N_+$ e $N_-$ soddisfano $N = N_+ + N_-$. Inoltre noti $E, N$ possiamo esprimere $N_+$ e $N_-$ in funzione di essi, infatti 
\begin{align}
    \begin{cases}
        N_+ &= \frac{1}{2}\left(N + \frac{E}{\varepsilon}\right) \\
        N_- &= \frac{1}{2}\left(N - \frac{E}{\varepsilon}\right) 
    \end{cases} \tag{2.6.2}
\end{align}
e quindi esprimere $\Gamma$ in funzione di $N, E$ 
\begin{align}
    \Gamma(E) = \frac{N!}{\left(\frac{1}{2}\left(N + \frac{E}{\varepsilon}\right)\right)!\left(\frac{1}{2}\left(N - \frac{E}{\varepsilon}\right)\right)!} \tag{2.6.3}
\end{align}
Usando la formula di Stirling al limite termidinamico si ha 
\begin{align}
    \Gamma(E) \sim \frac{N^{N + \frac{1}{2}}}{\left(\frac{1}{2}\left(N + \frac{E}{\varepsilon}\right)\right)^{\frac{1}{2}\left(N + \frac{E}{\varepsilon}\right) + \frac{1}{2}}\left(\frac{1}{2}\left(N - \frac{E}{\varepsilon}\right)\right)^{\frac{1}{2}\left(N - \frac{E}{\varepsilon}\right) + \frac{1}{2}}} \tag{2.6.4}
\end{align}
Segue che 
\begin{align}
    S(E) &\sim o(\log N) + k\left(N\log N - \frac{1}{2}\left(N + \frac{E}{\varepsilon}\right)\log\frac{1}{2}\left(N + \frac{E}{\varepsilon}\right) +\right. \notag \\
    &-\left.\frac{1}{2}\left(N - \frac{E}{\varepsilon}\right)\log\frac{1}{2}\left(N - \frac{E}{\varepsilon}\right) - N\log\frac{1}{2}\right) \tag{2.6.5}
\end{align}
dove 
\begin{align}
    o(\log N) &= \frac{k}{2}\log N - \frac{k}{2}\log\frac{1}{2}\left(N + \frac{E}{\varepsilon}\right) - \frac{k}{2}\log\frac{1}{2}\left(N - \frac{E}{\varepsilon}\right) \notag \\
    &= \frac{k}{2}\log\frac{N}{\frac{1}{4}\left(N^2 - \frac{E^2}{\varepsilon^2}\right)} \notag \\
    &= \frac{k}{2}\log\frac{4\varepsilon^2}{kT^2c_V} \tag{2.6.6}
\end{align}
Studiando lo stesso sistema nel contesto dell'ensemble canonico, possiamo determinare l'energia libera $F$ (e quindi l'entropia canonica) passando per la funzione di partizione canonica
\begin{align}
    Q_N &= \sum_n e^{-\beta E_n} \notag \\
    &= \sum_{N_+ + N_- = N}\frac{N!}{N_+!N_-!}e^{-\beta(N_+\varepsilon - N_-\varepsilon)} \notag \\
    &= \sum_{k = 0}^N\frac{N!}{k!(N - k)!}e^{-\beta\varepsilon k}e^{\beta\varepsilon(N - k)} \notag \\
    &= \left(e^{\beta E} + e^{-\beta E}\right)^N = Q^N_1 \tag{2.6.7}
\end{align}
dove si è riconosciuta l'espressione del \textit{binomio di Newton}. A questo punto possiamo determinare l'espressione dell'energia libera 
\begin{align}
    F &= -NkT\log(e^{\beta E} + e^{-\beta E}) \notag \\
    &= -NkT\log(2\cosh \beta E) \tag{2.6.8}
\end{align}
Segue che 
\begin{align}
    U &= -NE\frac{e^{\beta E} - e^{-\beta E}}{e^{\beta E} + e^{-\beta E}} \notag \\
    &= -NE\tanh \beta E \tag{2.6.9}
\end{align}
Per definizione, l'entropia canonica del sistema è data da
\begin{align}
    S_{can} &= -\left(\frac{\partial F}{\partial T}\right)_{N, V} \notag \\
    &= \frac{\partial}{\partial T}\left(NkT\log\left(2\cosh \frac{E}{kT}\right)\right) \notag \\
    &= Nk\log\left(2\cosh \frac{E}{kT}\right) - \frac{NE}{T}\tanh\frac{E}{kT} \notag \\
    &= Nk\log\left(2\cosh \beta E\right) - N\beta kE\tanh \beta E \tag{2.6.10}
\end{align}

\newpage
\subsection{Sistema di $N$ oscillatori armonici quantisitici}
Consideriamo un sistema composto da $N$ oscillatori armonici quantistici. Dato che lo spettro energetico è discreto si ha 
\begin{align}
    Q_N = \sum_{states} e^{-\beta E_{state}} \tag{2.7.1}
\end{align}
Anche in questo caso le particelle sono indipendenti e distinguibili, quindi $Q_N = Q^N_1$, dove 
\begin{align}
    Q_1 &= \sum_n e^{-\beta E_n} \notag \\
    &= \sum_n e^{-\beta\hbar\omega(n + 1/2)} \notag \\
    &= e^{-\beta\frac{\hbar\omega}{2}}\sum_n e^{-\beta\hbar\omega n} \tag{2.7.2}
\end{align}
con $E_n = \hbar\omega\left(n + \frac{1}{2}\right)$. Poniamo $x = e^{-\beta\hbar\omega}$ e notiamo che $|x| \leq 1$, infatti $\beta\hbar\omega > 0$, pertanto possiamo ricondurre $(2.7.2)$ alla serie geometrica di ragione $x$, segue che 
\begin{align}
    Q_1 &= e^{-\beta\frac{\hbar\omega}{2}}\frac{1}{1 - x} \notag \\
    &= \frac{e^{-\beta\frac{\hbar\omega}{2}}}{1 - e^{-\beta\hbar\omega}} \tag{2.7.3}
\end{align}
Allora 
\begin{align}
    Q_N = e^{-N\beta\frac{\hbar\omega}{2}}\left(1 - e^{-\beta\hbar\omega}\right)^{-N} \tag{2.7.4}
\end{align}
Nota la funzione di partizione canonica possiamo determinare l'espressione dell'energia libera 
\begin{align}
    F &= N\frac{\hbar\omega}{2} + NkT\log\left(1 - e^{-\beta\hbar\omega}\right) \tag{2.7.5}
\end{align}
Mentre per l'energia interna si ha 
\begin{align}
    U &= -\frac{\partial}{\partial\beta}\log Q_N \notag \\
    &= -\frac{\partial}{\partial\beta}\left(-N\beta\frac{\hbar\omega}{2} - N\log\left(1 - e^{-\beta\hbar\omega}\right)\right) \notag \\
    &= N\hbar\omega\left(\frac{1}{2} - \frac{e^{-\beta\hbar\omega}}{1 - e^{-\beta\hbar\omega}}\right) \tag{2.7.6}
\end{align}
Al limite per cui vale $\frac{\hbar\omega}{kT} << 1$ si ha 
\begin{align}
    U &\sim N\cancel{\frac{\hbar\omega}{2}} + N\hbar\omega\frac{1}{\frac{\hbar\omega}{kT}} \notag \\
    &= NkT \tag{2.7.7}
\end{align}
Il caso di $N$ oscillatori armonici quantistici tridimensionali è equivalente a avere $3N$ oscillatori armonici quantistici unidimensionali, pertanto 
\begin{align}
    Q_N = e^{-\frac{3}{2}\beta\hbar\omega}\left(1 - e^{-\beta\hbar\omega}\right)^{-3N} \tag{2.7.8}
\end{align}
e 
\begin{align}
    U = 3NkT \tag{2.7.9}
\end{align}
come previsto dal \textit{teorema di equipartizione dell'energia}.

\newpage
\subsection{Entropia canonica}
Con un opportuno cambio di variabili si può esprimere la funzione di partizione canonica come integrale sui livelli energetici del sistema (che supponiamo continui), ovvero
\begin{align}
    Q_N = \frac{1}{h^{3N}}\int d\mathbf{p}d\mathbf{q} \, e^{-\beta\mathcal{H}} = \int_0^\infty dE \, \Gamma(E)e^{-\beta E} \tag{2.8.1}
\end{align}
Essenzialmente stiamo integrando su tutte le ipersuperficie a energia fissata nello spazio delle fasi. Per definizione 
\begin{align}
    \Gamma(E) = \exp\left(-\frac{S_{mic}(E)}{k}\right) \tag{2.8.2}
\end{align}
Sostituendo $(2.8.2)$ in $(2.().1)$ si ha 
\begin{align}
    Q_N &= \int_0^\infty dE \, \exp\left(-\frac{S_{mic}(E)}{k} - \beta E\right) \notag \\
    &= \int_0^\infty dE \, \exp\left(-\beta\left(TS_{mic}(E) + E\right)\right) \tag{2.8.3}
\end{align}
Notiamo che la funzione integranda è prodotto di una funzione crescente e una funzione decrescente (dell'energia) ed ha quindi un punto di massimo globale. \\
A questo punto, sotto la stessa ipotesi dell'esempio del gas libero, ha senso approssimare l'esponente attorno al punto di massimo. Si ha 
\begin{align}
    \frac{\partial}{\partial E}(TS_{mic}(E) - E) = 0 \notag
\end{align}
per $E = \overline{E}$ tale che
\begin{align}
    T\left(\frac{\partial S_{mic}}{\partial E}\right)_{E = \overline{E}} = 1 \notag
\end{align}
Attorno ad $\overline{E}$ si ha 
\begin{align}
    TS_{mic}(E) - E \sim TS_{mic}(\overline{E}) - \overline{E} + \frac{1}{2}(E - \overline{E})^2T\left(\frac{\partial^2S_{mic}}{\partial E^2}\right)_{E = \overline{E}} \tag{2.8.4}
\end{align}
Ponendo $\eta = E - \overline{E}$ si ottiene 
\begin{align}
    Q_N \sim \frac{1}{h^{3N}}\exp\left(\beta(TS_{mic}(\overline{E}) - \overline{E})\right)\int_0^\infty d\eta \, \exp\left(\frac{1}{2k}\left(\frac{\partial^2S_{mic}}{\partial E^2}\right)_{E = \overline{E}}\eta^2\right) \tag{2.8.5}
\end{align}
Notiamo che la derivata seconda di $S_{mic}$ calcolata per $E = \overline{E}$ è un termine negativo, infatti $\overline{E}$ è un punto di massimo per $S_{mic}$, segue che l'integrale in $(2.8.5)$ è gaussiano, quindi 
\begin{align}
    Q_N \sim \frac{1}{2h^{3N}}\exp\left(\beta(TS_{mic}(\overline{E}) - \overline{E})\right)\left(-2\pi k\left(\frac{\partial^2S_{mic}}{\partial E^2}\right)^{-1}_{E = \overline{E}}\right)^{1/2} \tag{2.8.6}
\end{align}
Concludiamo che l'entropia canonica è maggiore dell'entropia microcanonica (a causa delle fluttazioni) e differisce da quest'ultima per un termine che va come $\log N$. In particolare, usando la definizione statistica dell'energia libera, si ottiene 
\begin{align}
    U - TS_{can} \sim \overline{E} - TS_{mic} - \frac{kT}{2}\log\left(2\pi kT^2\frac{\partial \overline{E}}{\partial T}\right) \tag{2.8.7}
\end{align}
e, dato che $U \sim \overline{E}$, si ottiene 
\begin{align}
    S_{can} \sim S_{mic} + \frac{kT}{2}\log\left(2\pi kT^2c_V\right) \tag{2.8.8}
\end{align}

\newpage
\section{Ensemble grancanonico}
Siamo interessati a descrivere microscopicamente un sistema a contatto con un reservoir, dunque termalizzato ad una temperatura $T$, con il quale scambia energia ed è in grado di interagire chimicamente. Segue che le grandezze termodinamiche $E, N$ diventano ora variabili stocastiche del sistema. \\
Sistemi che soddisfano tal condizioni fanno parte dell'\textit{ensemble grancanonico}.
\subsection{Legge di stato}
Consideriamo un sistema $X$ a contatto con un reservoir termalizzato e suddividiamo $X$ in due sottosistemi $X_1, X_2$ caratterizzati dalle grandezze termodinamiche $N_1, V_1$ e $N_2, V_2$, che soddisfano 
\begin{align}
    N &= N_1 + N_2 \notag \\
    V &= V_1 + V_2 \notag
\end{align}
Al limite in cui $V$ è molto grande il contributo di interazione tra $X_1$ e $X_2$, che è pressoché limitato al bordo di $X_1$, è trascurabile rispetto al contributo dovuto alle interazioni interne, ovvero i due sistemi sono debolmente interagenti e quindi si ha 
\begin{align}
    \mathcal{H}_N(\mathbf{p}, \mathbf{q}) = \mathcal{H}_{N_1}(\mathbf{p}_1, \mathbf{q}_1) + \mathcal{H}_{N_2}(\mathbf{p}_2, \mathbf{q}_2) \notag
\end{align}
Per definizione
\begin{align}
    Q_N &= \frac{1}{h^{3N}}\int d\mathbf{p}d\mathbf{q} \, e^{-\beta\mathcal{H}_N} \notag
\end{align}
Cerchiamo un modo per scrivere la funzione di partizione canonica come segue (costruzione marginale (?))
\begin{align}
    Q_N = \sum_{N_1}\frac{1}{h^{3N_1}}\int d\mathbf{p}_1d\mathbf{q}_1 \, (\dots) \tag{3.1.1}
\end{align}
Poiché ora stiamo considerando $N_1$ variabile, ovvero supponiamo che il sottosistema $X_1$ possa interagire chimicamente con $X_2$, dobbiamo considerare tutte le configurazioni possibili che soddisfano il vincolo $N = N_1 + N_2$. \\
Se $N_1$ e $N_2$ fossero fissate, sotto l'ipotesi di sistemi debolmente interagenti, si avrebbe
\begin{align}
    Q_N &= Q_{N_1}Q_{N_2} \notag \\
    &= \left(\frac{1}{\hbar^{3N_1}N_1!}\int d\mathbf{p}^{(1)}d\mathbf{q}^{(1)} \, e^{-\beta\mathcal{H}_{N_1}}\right)\left(\frac{1}{\hbar^{3N_2}N_2!}\int d\mathbf{p}^{(2)}d\mathbf{q}^{(2)} \, e^{-\beta\mathcal{H}_{N_2}}\right) \tag{3.1.2}
\end{align}
Una volta introdotta la possibilità che i sistemi scambino particelle dobbiamo considerare tutte le configurazioni che soddisfano $N = N_1 + N_2$, pertanto si ottiene
\begin{align}
    Q_N &= \sum_{N_1 = 0}^N \left(\frac{1}{\hbar^{3N_1}N_1!}\int d\mathbf{p}^{(1)}d\mathbf{q}^{(1)} \, e^{-\beta\mathcal{H}_{N_1}}\right)\left(\frac{1}{\hbar^{3N_2}N_2!}\int d\mathbf{p}^{(2)}d\mathbf{q}^{(2)} \, e^{-\beta\mathcal{H}_{N_2}}\right) \notag \\
    &= \sum_{N_1 = 0}^N \frac{1}{N_1!(N - N_1)!}Q_{N_1}Q_{N - N_1} \tag{3.1.3}
\end{align}
Concludiamo che la densità di probabilità associata al sottosistema $X_1$ è data da
\begin{align}
    \rho_{N_1}(\mathbf{p}_1, \mathbf{q}_1) &= \frac{Q_{N_2}}{Q_N}\frac{e^{-\beta\mathcal{H}_{N_1}}}{N_1!h^{3N_1}} \notag \\
    &= e^{-\beta(F_{N_2}(T, V_2) - F_N(T, V))}\frac{e^{-\beta\mathcal{H}_{N_1}}}{N_1!h^{3N_1}} \tag{3.1.4}
\end{align}
Notiamo che $\rho$ è correttamente normalizzata sullo spazio delle fasi associato a $X_1$, in particolare 
\begin{align}
    1 = \sum_{N_1 = 0}\int d\mathbf{p}_1d\mathbf{q}_1 \, \rho_{N_1}(\mathbf{p}_1, \mathbf{q}_1) \tag{3.1.5}
\end{align}
Supponiamo ora che $N_1 << N$ e $V_1 << V$, si ha 
\begin{align}
    F_{N_2}(T, V_2) - F_N(T, V) &\sim -\underbrace{\left(\frac{\partial F}{\partial N}\right)_{T, V}}_{\mu}N_1 - \underbrace{\left(\frac{\partial F}{\partial V}\right)_{N, T}}_{-P}V_1 \notag \\
    &= -\mu N_1 + PV_1 \tag{3.1.6}
\end{align}
sostituendo in $(3.1.4)$ si ottiene 
\begin{align}
    \rho_{N_1}(\mathbf{p}_1, \mathbf{q}_1) \sim e^{-\beta(PV_1 -\mu N_1)}\frac{e^{-\beta \mathcal{H}_{N_1}}}{N_1!h^{3N_1}} \tag{3.1.7}
\end{align}
A questo punto possiamo abbandonare il pedice che caratterizza il sottosistema $X_1$ e lo trattiamo come il sistema generale. Per un insieme di particelle indipendenti e indistinguibili nel contesto dell'ensemble grancanonico si ha
\begin{align}
    \rho_N(\mathbf{p}, \mathbf{q}) = e^{-\beta(PV + \mathcal{H}_N)}\frac{Z^N}{N!h^{3N}} \tag{3.1.8}
\end{align}
dove si è posto $Z = e^{\beta\mu}$, detta \textit{fugacità}. Sostituendo l'espressione $(3.1.7)$ in $(3.1.5)$ otteniamo 
\begin{align}
    1 &= \sum_{N = 0}^{\overbrace{\infty}^{\text{lim. termodin.}}}e^{-\beta PV}Z^N\frac{1}{N!h^{3N}}\int d\mathbf{p}d\mathbf{q} \, e^{-\beta\mathcal{H}_N} \notag \\
    &= \sum_{N = 0}^\infty e^{-\beta PV}Z^NQ_N \tag{3.1.9}
\end{align}
ovvero 
\begin{align}
    e^{\beta PV} &= \sum_{N = 0}^\infty Z^NQ_N \tag{3.1.10}
\end{align}
Poniamo 
\begin{align}
    \tilde{Q}(T, V, Z) = \sum_{N = 0}^\infty Z^NQ_N \tag{3.1.11}
\end{align}
detta \textit{funzione di partizione grancanonica} (costante che determina la normalizzazione della funzione distribuzione di probabilità $\rho$ associata all'ensemble grancanonico). \\
Per particelle indipendenti e indistinguibili sappiamo che vale $Q_N = \frac{Q^N_1}{N!}$. Sostituendo in $(3.1.10)$ si ottiene 
\begin{align}
    e^{\beta PV} = \sum_{N = 0}^\infty \frac{(ZQ_1)^N}{N!} = e^{ZQ_1} \tag{3.1.12}
\end{align}
e quindi otteniamo 
\begin{align}
    PV = ZQ_1kT \tag{3.1.13}
\end{align}
Nel contesto dell'ensemble grancanonico, il numero di particelle è una variabile stocastica, a causa del potenziale chimico $\mu$. Tuttavia, come già visto nel caso dell'ensemble canonico, è ragionevole aspettarsi che $ZQ_1 = \braket{N}$. \\ \\
Notiamo che
\begin{align}
    \braket{N} &= \sum_{N = 0}^\infty \frac{1}{\tilde{Q}}\frac{Z^N}{N!h^{3N}}\int d\mathbf{p}d\mathbf{q} \, Ne^{-\beta \mathcal{H}_N} \notag \\
    &= \frac{1}{\tilde{Q}}\sum_{N = 0}^\infty NZ^NQ_N \notag \\
    &= Z\left(\frac{\partial}{\partial Z}\log \tilde{Q}\right)_{T, V} \tag{3.1.14}
\end{align}
dove nel secondo passaggio si è usata l'identità $(3.1.10)$. Pertanto, sempre nel contesto di particelle indipendenti e indistinguibili, si ha 
\begin{align}
    \braket{N} &= Z\frac{\partial}{\partial Z}\left(\log \tilde{Q}\right) \notag \\
    &= Z\frac{\partial}{\partial Z}\left(\log \sum_{N = 0}^\infty Z^NQ_N\right) \notag \\
    &= Z\frac{\partial}{\partial Z}\left(\log \sum_{N = 0}^\infty \frac{(ZQ_1)^N}{N!}\right) \notag \\
    &= Z\frac{\partial}{\partial Z}ZQ_1 = ZQ_1 \tag{3.1.15}
\end{align}
dove si è usato $(3.1.12)$ nel quarto passaggio. Quindi, come ci si aspettava, si ha 
\begin{align}
    PV = \braket{N}kT \tag{3.1.16}
\end{align}
Inoltre, grazie alla $(3.1.15)$, siamo in grado di ottenere un'espressione per il potenziale chimico a cui è soggetto il sistema passando per la definizione di fugacità, in particolare si ha 
\begin{align}
    \mu = -kT\log \frac{Q}{\braket{N}} \tag{3.1.17}
\end{align}

\newpage
\subsection{Fluttuazioni del numero medio di particelle (spettro discreto)}
Poniamoci sempre nel contesto di particelle indipendenti e indistinguibili. Come già visto la funzione di partizione canonica si fattorizza. Supponiamo che $Q_1 = Vf(T)$, con $f$ una qualsiasi funzione della temperatura sufficientemente regolare. \\ \\
Nel caso di spettro discreto si ha 
\begin{align}
    \tilde{Q} &= \sum_{N = 0}^\infty Z^NQ_N = \sum_{N, E_N} Z^Nd(E_N)e^{-\beta E_N} \tag{3.2.1}
\end{align}
dove $d(E_N)$ rappresenta la degenerazione del livello energetico $E_N$. Possiamo quindi definire una distribuzione di probabilità $P$ ponendo 
\begin{align}
    P(N, E_N) = \frac{Z^Nd(E_N)e^{-\beta E_N}}{\sum_{N, E_N}Z^Nd(E_N)e^{-\beta E_N}} \tag{3.2.2}
\end{align}
ben normalizzata. A questo punto, per costruzione, si ha 
\begin{align}
    U = \sum_{N, E_N} E_NP(N, E_N) = -\left(\frac{\partial}{\partial\beta}\log \tilde{Q}\right) _{V, Z}\notag
\end{align}
Per $(3.1.12)$ si ha che 
\begin{align}
    -\left(\frac{\partial}{\partial\beta}\log \tilde{Q}\right)_{V, Z} = -ZV\frac{\partial}{\partial\beta}f(T) = ZVf'(T)kT^2 \tag{3.2.3}
\end{align}
dove si è usata l'ipotesi $\tilde{Q} = Vf(T)$. A questo punto possiamo calcolare il calore specifico a volume costante $c_V$ usando la definizione termodinamica 
\begin{align}
    c_V &= \left(\frac{\partial}{\partial T}\braket{E}\right)_{V, Z} \notag \\
    &= \left(\frac{\partial}{\partial T}\frac{f'(T)}{f(T)}\braket{N}kT^2\right)_{V, Z} \notag \\
    &= \braket{N}kT\frac{2f(T)f'(T) + Tf(T)f''(T) - T(f'(T))^2}{(f(T))^2} \tag{3.2.4}
\end{align}
dove si è usato $\braket{N} = ZVf(T)$ e quindi $Z = \frac{\braket{N}}{Vf(T)}$. Questa espressione tornerà molto utile quando la dipendenza di $f$ da $T$ è della forma $T^\alpha$ per qualche $\alpha \in \mathbf{R}$. \\ \\
Per $(3.1.12)$ si ha che 
\begin{align}
    P(N, E_N) &= \frac{Z^Nd(E_N)e^{-\beta E_N}}{\sum_{N, E_N}Z^Nd(E_N)e^{-\beta E_N}} \notag \\
    &= \frac{Z^NQ_N}{\sum_N Z^NQ_N} \notag \\
    &= \frac{Z^NQ^N_1}{N!}e^{-ZQ_1} \notag \\
    &= \frac{\braket{N}}{N!}e^{-\braket{N}} \tag{3.2.5}
\end{align}
ovvero $P(N, E_N)$ è una distribuzione poissoniana. Nell'appendice si dimostra che per $\braket{N} \rightarrow \infty$ la distribuzione poissoniana si comporta come una distribuzione gaussiana. \\ \\
A questo punto studiamo come si comportano le fluttuazioni del numero di particelle. Per definizione 
\begin{align}
    \braket{N^2} &= \frac{\sum_{N}N^2Z^NQ_N}{\sum_{N}Z^NQ_N} \notag
\end{align}
e 
\begin{align}
    \braket{(\Delta N)^2} &= \braket{N^2} - \braket{N}^2 \notag \\
    &= \frac{\sum_{N}N^2Z^NQ_N}{\sum_{N}Z^NQ_N} - \left(\frac{\sum_{N}NZ^NQ_N}{\sum_{N}Z^NQ_N}\right)^2 \tag{3.2.6}
\end{align}
Verifichiamo che vale la seguente relazione 
\begin{align}
    \braket{(\Delta N)^2} = Z\left(\frac{\partial}{\partial Z}\braket{N}\right)_{V, T} \tag{3.2.7}
\end{align}
Si ha che 
\begin{align}
    Z\left(\frac{\partial}{\partial Z}\braket{N}\right)_{V, T} &= Z\left(\frac{\partial}{\partial Z}\frac{\sum_{N}NZ^NQ_N}{\sum_{N}Z^NQ_N}\right)_{V, T} \notag \\
    &= Z\left(\frac{\sum_{N}N^2Z^{N - 1}Q_N}{\sum_{N}Z^NQ_N} - \frac{\sum_{N}NZ^{N - 1}Q_N\sum_{N}NZ^NQ_N}{\left(\sum_N Z^NQ_N\right)^2}\right) \notag \\
    &= \frac{\sum_{N}N^2Z^NQ_N}{\sum_{N}Z^NQ_N} - \left(\frac{\sum_{N}NZ^NQ_N}{\sum_{N}Z^NQ_N}\right)^2 \notag \\
    &= \braket{N^2} - \braket{N}^2 \tag{3.2.8}
\end{align}
Pertanto $(3.2.7)$ è dimostrata.


\newpage
\subsection{Particelle libere in una scatola}
Consideriamo un sistema di $N$ particelle libere, indipendenti e indistinguibili in una scatola di volume $V$. In questo caso sappiamo che vale $Q_N = \frac{Q^N_1}{N!}$, inoltre 
\begin{align}
    Q_1 = cVT^{3/2} \tag{3.3.1}
\end{align}
dove $c$ è un termine costante. Come fatto per $\braket{N}$, usando la definizione di media per l'ensemble grancanonico, possiamo scrivere 
\begin{align}
    U = -\left(\frac{\partial}{\partial \beta}\log\tilde{Q}\right)_{V, Z} \tag{3.3.2}
\end{align}
Allora si ottiene 
\begin{align}
    U &= -\left(\frac{\partial}{\partial \beta}ZQ_1(T, V)\right)_{V, Z} \notag \\
    &= -Z\frac{\partial T}{\partial \beta} \left(\frac{\partial}{\partial T}Q_1(T, V)\right)_{V, Z} \notag \\
    &= \frac{3}{2}\overbrace{cZVT^{3/2}}^{\braket{N}}kT = \frac{3}{2}\braket{N}kT \tag{3.3.3}
\end{align}
Mentre ricordiamo che per l'ensemble canonico si ha $U_{can} = \frac{3}{2}NkT$ con $N$ costante del sistema.

\newpage
\section{Meccanica Statistica Quantistica}
Nel caso di sistemi quantistici è necessario tenere conto di aspetti che abbiamo fino ad ora trascurato, come l'indistinguibilità quantistica delle particelle, che porta a funzioni d'onda con simmetria ben definita, e la sovrapposizione di stati.

\subsection{Indistinguibilità quantistica}
Consideriamo un sistema composto da due particelle indipendenti, ognuna delle quali può avere un energia che è multiplo intero positivo di una quantità fissata $\varepsilon > 0$. L'energia totale del sistema è data da
\begin{align}
    E_{tot} = m\varepsilon + n\varepsilon \qquad m, n \in \mathbf{N} \tag{4.1.1}
\end{align}
Possiamo calcolare la funzione di partizione canonica associata al sistema, in particolare
\begin{align}
    Q_2 = \sum_{states} e^{-\beta E_{state}} = \left(\sum_{m = 0}^\infty e^{-\beta m \varepsilon}\right)\left(\sum_{n = 0}^\infty e^{-\beta n \varepsilon}\right) = Q^2_1 \tag{4.1.2}
\end{align}
Dato che si tratta di particelle indipendenti, la funzione di partizione canonica si scompone nel prodotto di funzioni di partizione canonica di singola particella. \\
Se teniamo conto anche dell'indistinguibilità quantistica e delle proprietà di simmetria dovute alla natura delle particelle il risultato cambia. \\
Supponiamo che si tratti di due fermioni. In tal caso la funzione d'onda complessiva deve essere asimmetrica, quindi le due particelle non possono trovarsi nello stesso stato, ovvero $n = m$ è vietato. Inoltre, $(m, n)$ e $(n, m)$ identificano lo stesso stato. Pertanto otteniamo
\begin{align}
    Q_2 = \frac{1}{2}\sum_{m, n = 0}^\infty e^{-\beta(m\varepsilon + n\varepsilon)} - \frac{1}{2}\sum_{m = n}e^{-\beta(m\varepsilon + n\varepsilon)} = \frac{1}{2}Q^2_1(\beta) - \frac{1}{2}Q_1(2\beta) \tag{4.1.3}
\end{align}
Mentre, se supponiamo che si tratti di due bosoni, la funzione d'onda complessiva deve essere simmetrica, quindi le particelle possono occupare lo stesso stato ($m = n$ ammesso). Come prima stati $(m, n)$ e $(n, m)$ identificano la stessa configurazione, dato che le particelle sono indistinguibili. Pertanto 
\begin{align}
    Q_2 = \frac{1}{2}\sum_{m, n = 0}^\infty e^{-\beta(m\varepsilon + n\varepsilon)} + \frac{1}{2}\sum_{m = n}e^{-\beta(m\varepsilon + n\varepsilon)} = \frac{1}{2}Q^2_1(\beta) + \frac{1}{2}Q_1(2\beta) \tag{4.1.4}
\end{align}
Consideriamo un sistema simile, composto da due particelle che possono occupare due stati differenti, uno a energia $0$ e l'altro a energia $\varepsilon$. La funzione di partizione canonica è data da
\begin{align}
    Q_2 = 1 + 2e^{-\beta\varepsilon} + e^{-2\beta\varepsilon} = (1 + e^{-\beta\varepsilon})^2 = Q^2_1 \tag{4.1.5}
\end{align}
Se supponiamo che le particelle siano indistinguibili e moltiplichiamo per il fattore di Gibbs, si ottiene
\begin{align}
    Q_2 = \frac{1}{2} + e^{-\beta\varepsilon} + \frac{1}{2}e^{-2\beta\varepsilon} \tag{4.1.6}
\end{align}
Tuttavia, non è corretto assegnare un fattore $\frac{1}{2}$ ad ogni microstato. \\
Se consideriamo le particelle come indistinguibili sin dall'inizio i microstati distinti diventano tre, di cui solo uno è accessibile nel caso di fermioni. In particolare, per fermioni si ottiene
\begin{align}
    Q_2 = e^{-\beta\varepsilon} \tag{4.1.7}
\end{align}
che combacia con il risultato ottenuto in $(4.1.6)$ se non consideriamo gli stati non accessibili. Mentre, nel caso di bosoni si ottiene 
\begin{align}
    Q_2 = 1 + e^{-\beta\varepsilon} + e^{-2\beta\varepsilon} \tag{4.1.8}
\end{align}
Notiamo che $(4.1.8)$ e $(4.1.6)$ sono differenti. Infatti, il fattore di Gibbs fallisce quando c'è un accumulo di particelle in un singolo stato. Più precisamente, quando la frazione di particelle in un singolo stato diventa non infinitesima, il fattore di Gibbs fallisce. \\
Questo accade perché il fattore di Gibbs non selezione solamente i microstati "patologici", ma associa un peso statistico identico ad ogni microstato. In $(4.1.8)$ solamente i microstati che contruibiscono con il termine $e^{-\beta\varepsilon}$ creano problemi, ovvero gli stati $(\varepsilon, 0)$ e $(0, \varepsilon)$ che sono identici per particelle indistinguibili. Mentre, i microstati $(0, 0)$ e $(\varepsilon, \varepsilon)$ non vengono contati più volte nel calcolo della funzione di partizione canonica, pertanto non necessitano del peso statistico attribuito dal fattore di Gibbs.

\newpage
\subsection{Concetto di ensemble in meccanica quantistica}
Consideriamo un sistema quantistico isolato a più corpi descritto dall'hamiltoniano $\hat{\mathcal{H}}$. Un qualsiasi stato quantistico puro (che presenti solamente fluttazioni quantistiche) $\ket{\psi(t)}$ può essere espresso come combinazione lineare di autostati $\{\ket{n}\}_n$ dell'hamiltoniniano, ovvero 
\begin{align}
    \ket{\psi(t)} = \sum_n c_ne^{-i\frac{E_n}{\hbar}t}\ket{n} \tag{4.2.1}
\end{align}
dove $\sum_n |c_n|^2 = 1$, che garantisce la normalizzazione del ket di stato per ogni istante di tempo $t$. La condizione iniziale $\ket{\psi(0)}$ è univocamente identificata dalla scelta dei coefficienti $c_n$
\begin{align}
    \ket{\psi(0)} = \sum_n c_n\ket{n} \tag{4.2.2}
\end{align}
Consideriamo un'osservabile $\mathcal{O}$ associata al sistema in esame. Per definizione, il valore medio dell'osservabile $\mathcal{O}$ associata allo stato $\ket{\psi(t)}$ è data da 
\begin{align}
    \braket{\mathcal{O}}_t = \braket{\psi(t)|\mathcal{O}|\psi(t)} = \sum_n |c_n|^2\braket{n|\mathcal{O}|n} + \sum_{m \neq n}c_m^*c_ne^{-\frac{i}{\hbar}(E_m - E_n)t}\braket{m|\mathcal{O}|n} \tag{4.2.3}
\end{align}
Notiamo che la prima sommatoria in $(4.2.3)$ contribuisce con un termine indipendente dal tempo, mentre la seconda sommatoria contribuisce con un termine che dipende armonicamente dal tempo. \\
A questo punto, ci chiediamo se è possibile osservare equilibrio nel sistema. In generale sembrerebbe di no, infatti la somma sulla diagonale dipende fortemente dai coefficienti che definiscono la condizione iniziale del sistema, mentre la somma dei termini fuori diagonale ha dipendenza armonica dal tempo, che in generale non si annulla. \\ \\
Per studiare il comportamento all'equilibrio di sistemi quantistici a più corpi, consideriamo la media statistica su un ensemble costituito da $N$ copie del sistema in esame, ognuno con condizione iniziale $\ket{\psi_k(0)}$ differente. In tale schema, il valor medio sull'ensemble di un'osservabile $\mathcal{O}$ è dato da 
\begin{align}
    \braket{\mathcal{O}}_t = \frac{1}{N}\sum_{k = 0}^N \braket{\psi_k(t)|\mathcal{O}|\psi_k(t)} \tag{4.2.4}
\end{align}
Definiamo $p_k = \frac{n_k}{N}$, dove $n_k$ è il numero di elementi dell'ensemble che hanno come condizione iniziale $\ket{\psi_k(0)}$, notiamo che $\sum_k p_k = 1$. Allora 
\begin{align}
    \braket{\mathcal{O}}_t = \sum_k p_k\braket{\psi_k(t)|\mathcal{O}|\psi_k(t)} \tag{4.2.5}
\end{align}
Introducendo un set completo $\sum_\alpha \ket{\alpha}\bra{\alpha}$, dove $\{\ket{\alpha}\}_{\alpha}$ non sono necessariamente autostati dell'hamiltoniano $\hat{\mathcal{H}}$, si può riscrivere $(4.2.5)$ come segue 
\begin{align}
    \braket{\mathcal{O}}_t &= \sum_k\sum_\alpha p_k\braket{\psi_k(t)|\alpha}\braket{\alpha|\mathcal{O}|\psi_k(t)} \notag \\
    &= \sum_\alpha \braket{\alpha|\left(\sum_k p_k \ket{\psi_k(t)}\bra{\psi_k(t)}\right)\mathcal{O}|\alpha} \tag{4.2.6}
\end{align}
Definiamo
\begin{align}
    \hat{\rho}(t) = \sum_k p_k\ket{\psi_k(t)}\bra{\psi_k(t)} \tag{4.2.7}
\end{align}
detta \textit{matrice densità}. Allora $(4.2.6)$ diventa 
\begin{align}
    \braket{O}_t = \sum_\alpha \braket{\alpha|\hat{\rho}(t)\mathcal{O}|\alpha} = \mathrm{Tr}(\hat{\rho}(t)\mathcal{O}) \tag{4.2.8}
\end{align}
Notiamo subito che la media sull'ensemble dell'osservabile $\mathcal{O}$ è indipendente dalla base $\{\ket{\alpha}\}_\alpha$ scelta, infatti la traccia è invariante per cambiamento di base. \\ \\
Vediamo come evolve nel tempo la matrice di densità
\begin{align}
    i\hbar\frac{\partial}{\partial t}\hat{\rho} &= \sum_k p_k(\hat{\mathcal{H}}\ket{\psi(t)}\bra{\psi(t)} - \ket{\psi(t)}\bra{\psi(t)}\hat{\mathcal{H}}) \notag \\
    &= \hat{\mathcal{H}}\hat{\rho} - \hat{\rho}\hat{\mathcal{H}} \notag \\
    &= [\hat{\mathcal{H}}, \hat{\rho}] \tag{4.2.9}
\end{align}
dove si è usato il fatto che lo stato quantistico in esame soddisfa l'equazione di Schrodinger dipendente dal tempo. Ricordiamo che per sistemi all'equilibrio deve valere
\begin{align}
    \frac{\partial}{\partial t}\hat{\rho} = 0 \tag{4.2.10}
\end{align}
Pertanto, confrontando $(4.2.10)$ con $(4.2.9)$, concludiamo che, se $\hat{\rho}$ è la matrice densità associata a un sistema in equilibrio termodinamico, allora è possibile diagonalizzarla simultaneamente all'hamiltoniano. Quindi si ottiene
\begin{align}
    \hat{\rho} = \sum_n \rho(E_n)\ket{n}\bra{n} \tag{4.2.11}
\end{align}
dove $\rho$ è una funzione scalare che attribuisce dei pesi statistici a ogni stato. In particolare, gli elementi di matrice della matrice densità sono
\begin{align}
    \hat{\rho}_{mn} = \rho(E_n)\delta_{mn} \tag{4.2.12}
\end{align}
L'elemento di diagonale $n$-esimo della matrice densità non è altro che una misura della probabilità di trovare il sistema nello stato $\ket{n}$. Naturalmente, la probabilità di trovarlo in tale stato dipende dall'energia $E_n$ a cui esso si trova. \\
Inoltre, la dipendenza di $\rho$ da $E_n$ cambia in base a quale ensemble stiamo considerando
\begin{align}
    \rho(E_n) &= \begin{cases}
        \frac{1}{\Gamma(E)} \qquad E &\leq E_n \leq E + \Delta \\
        0 \qquad &\text{altrimenti}
    \end{cases} \qquad &\text{microcanonico} \notag \\
    \rho(E_n) &= \frac{1}{Q_N(\beta)}e^{-\beta E_n} \qquad &\text{canonico} \notag \\
    \rho(E_n) &= \frac{1}{\tilde{Q}}e^{-\beta E_n - \beta\mu N_n} \qquad &\text{grancanonico} \notag 
\end{align}
Scegliendo di usare una rappresentazione differente dalla rappresentazione nella base dell'hamiltoniano per la matrice densità, diventa necessario introdurre un postulato che garantisce la compatibilità con $(4.2.11)$. Tale postulato è detto \textit{postulato delle fasi random}, in pratica afferma
\begin{align}
    \overline{e^{\frac{i}{\hbar}(\theta^k_m - \theta^k_n)}} = 0 \tag{4.2.13}
\end{align}
per ogni condizione iniziale $\ket{\psi_k(0)}$ e ogni $m \neq n$ (a meno di degenerazioni). In particolare, $(4.1.13)$ equivale a supporre che lo stato quantistico in esame sia dato da una sovrapposizione incoerenti di autostati, la quale ha contributo nullo una volta che viene mediato su tempi sufficientemente lunghi.

\newpage
\subsection{Gas ideali quantistici}
Consideriamo un sistema quantistico composto da $N$ particelle indipendenti. La funzione di partizione canonico associata a tale sistema è data da
\begin{align}
    Q_N = \sum_E e^{-\beta E} \tag{4.3.1}
\end{align}
dove $E$ è l'energia associata ad un determinato autostato. Possiamo esprimere $E$ in termini delle energie $\varepsilon$ dei singoli livelli occupati, ovvero 
\begin{align}
    E = \sum_{\varepsilon} n_\varepsilon\varepsilon \tag{4.3.2}
\end{align}
dove $n_\varepsilon$ è il numero di particelle che occupa lo stato ad energia $\varepsilon$. Ovviamente si ha 
\begin{align}
    \sum_\varepsilon n_\varepsilon = N \tag{4.3.3}
\end{align}
Pertanto, possiamo riscrivere la funzione di partizione canonica come segue
\begin{align}
    Q_N = \sum_{\{n_\varepsilon\}} g\{n_\varepsilon\}e^{-\beta\sum_\varepsilon n_\varepsilon\varepsilon} \tag{4.3.4}
\end{align}
dove $\{n_\varepsilon\}$ è l'insieme delle possibili configurazioni che soddisfano il vincolo $(4.3.3)$, mentre $g\{n_\varepsilon\}$ è il peso statistico assegnato ad una data configurazione. In particolare, il valore di $g$ per una data configurazione differesci nel caso di fermioni e bosoni. Per bosoni si ha
\begin{align}
    g\{n_\varepsilon\} = 1 \tag{4.3.5}
\end{align}
Mentre, nel caso di fermioni dobbiamo ricordare che uno stato non può essere occupato simultaneamente da più di un fermione, pertanto
\begin{align}
    g\{n_\varepsilon\} = \begin{cases}
        1 \qquad &\text{se la conf. è permessa} \\
        0 \qquad &\text{se la conf. è vietata}
    \end{cases} \tag{4.3.6}
\end{align}
Inoltre, è possibile riscrivere $(4.3.4)$ come 
\begin{align}
    Q_N = \sum_{\{n_\varepsilon\}}\prod_\varepsilon g\{n_\varepsilon\}\left(e^{-\beta\varepsilon}\right)^{n_\varepsilon} \tag{4.3.7}
\end{align}
A questo punto, dato che al limite termodinamico i tre ensemble sono identici, scegliamo l'ensemble grancanonico per studiare il comportamento all'equilibrio dei sistemi quantistici, dato che risulta essere il più semplice da trattare matematicamente. \\
La funzione di partizione grancanonica può essere scritta come
\begin{align}
    \tilde{Q} &= \sum_{N = 0}^\infty Z^NQ_N \notag \\
    &= \sum_{N = 0}^\infty\sum_{\{n_\varepsilon\}}\prod_\varepsilon \left(Ze^{-\beta\varepsilon}\right)^{n_\varepsilon} \notag \\
    &= \sum_{n_0, n_1, \dots}\left[\left(Ze^{-\beta\varepsilon_0}\right)^{n_0}\left(Ze^{-\beta\varepsilon_1}\right)^{n_1}\dots \right] \tag{4.3.8}
\end{align}
Ora diventa necessario distinguire il caso della statistica di Bose-Einstein con quello della statistica di Fermi-Dirac. \\
Nella prima situazione, ogni configurazione $\{n_\varepsilon\}$ che soddisfa il vincolo $(4.3.3)$ è permessa. Allora, nel limite termodinamico ($N \rightarrow \infty)$, si ottiene 
\begin{align}
    \tilde{Q}_{BE} = \prod_\varepsilon \frac{1}{1 - Ze^{-\beta\varepsilon}} \tag{4.3.9}
\end{align}
Mentre, nel caso di fermioni $n_\varepsilon$ può essere $0$ o $1$, pertanto
\begin{align}
    \tilde{Q}_{FD} = \prod_\varepsilon 1 + Ze^{-\beta\varepsilon} \tag{4.3.10}
\end{align}
Nota la funzione di partizione grancanonica, possiamo determinare le espressioni caratteristiche dell'ensemble grancanonico. In particolare
\begin{align}
    \frac{PV}{kT} &= \log \tilde{Q} = \begin{cases}
        -\sum_\varepsilon \log\left(1 - Ze^{-\beta\varepsilon}\right) &\qquad \text{BE}\\
        \sum_\varepsilon \log\left(1 + Ze^{-\beta\varepsilon}\right) &\qquad \text{FD}
    \end{cases} \tag{4.3.11} \\ \notag \\
    \mathcal{N} &= Z\frac{\partial}{\partial Z}\left(\log \tilde{Q}\right) = \begin{cases}
        \sum_\varepsilon \frac{1}{\frac{e^{\beta\varepsilon}}{Z} - 1} &\qquad \text{BE} \\
        \sum_\varepsilon \frac{1}{\frac{e^{\beta\varepsilon}}{Z} + 1} &\qquad \text{FD}
    \end{cases} \tag{4.3.12} \\ \notag \\
    U &= -\left(\frac{\partial}{\partial \beta \log\tilde{Q}}\right) = \begin{cases}
        \sum_\varepsilon \frac{\varepsilon}{\frac{e^{\beta\varepsilon}}{Z} - 1} &\qquad \text{BE} \\
        \sum_\varepsilon \frac{\varepsilon}{\frac{e^{\beta\varepsilon}}{Z} + 1} &\qquad \text{FD}
    \end{cases} \tag{4.3.13}
\end{align}
In ogni espressione è necessario moltiplicare per il termine $g_s = 2s + 1$ di degenerazione dovuta allo spin. \\ \\
Nel limite termodinamico gli stati energeetici sono densissimi, pertanto è possibile sostituire la sommatoria con un integrale nello spazio delle fasi
\begin{align}
    \sum_\varepsilon f(\varepsilon) \longrightarrow \int d\mathbf{p}d\mathbf{q} \, g(\mathbf{p}, \mathbf{q})f(\mathbf{p}, \mathbf{q}) \tag{4.3.14}
\end{align}
dove $g$ è la funzione densità degli stati e si è espressa la dipendenza di $f$ dall'energia come dipendenza da $(\mathbf{p}, \mathbf{q})$ (ovvero supponiamo che $\mathcal{H}$, hamiltoniano di singola particella, sia funzione solamente di coordinate e impulsi generalizzati). \\ \\
Osservato ciò, possiamo riscrivere le equazioni caratteristiche dell'ensemble grancanonico quantisico come segue
\begin{align}
    \frac{PV}{kT} &= (-1)^{g_s}g_s\int \frac{d\mathbf{p}d\mathbf{q}}{h^3} \, \log(1 + (-1)^{g_s}Ze^{-\beta\mathcal{H}(\mathbf{p}, \mathbf{q})}) \tag{4.3.15} \\ \notag \\
    \mathcal{N} &= (-1)^{g_s}g_s\int \frac{d\mathbf{p}d\mathbf{q}}{h^3} \, \frac{1}{Z^{-1}e^{\beta\mathcal{H}(\mathbf{p}, \mathbf{q})} + (-1)^{g_s}} \tag{4.3.16} \\ \notag \\
    U &= (-1)^{g_s}g_s\int \frac{d\mathbf{p}d\mathbf{q}}{h^3} \, \frac{\mathcal{H}(\mathbf{p}, \mathbf{q})}{Z^{-1}e^{\beta\mathcal{H}(\mathbf{p}, \mathbf{q})} + (-1)^{g_s}} \tag{4.3.17}
\end{align}
dove il valore di $g_s = 2s + 1$ dipende dalla natura delle particelle in esame (fermioni o bosoni).

\newpage
\subsection{Particelle libere}
Consideriamo un sistema composto da $N$ particelle indipendenti e indistinguibili (fermioni o bosoni), l'hamiltoniano di singola particelle che lo descrive è
\begin{align}
    \mathcal{H}(\mathbf{p}, \mathbf{q}) = \frac{\mathbf{p}^2}{2m} \notag
\end{align}
Nel contento dell'ensemble grancanonico quantistico si ha 
\begin{align}
    \frac{PV}{kT} = (-1)^{g_s}g_s\frac{V}{h^3}\int d\mathbf{p} \, \log(1 + (-1)^{g_s}Ze^{-\beta\frac{\mathbf{p}^2}{2m}}) \tag{4.4.1}
\end{align}
dove è stato subito calcolato il contributo dovuta all'integrazione in $d\mathbf{q}$. Passando in coordinare sferiche e integrando per parti $(4.4.1)$ si ottiene
\begin{align}
    \frac{PV}{kT} = &(-1)^{g_s}\frac{4\pi g_sV}{h^3}\left(\frac{1}{3}p^3\log(1 + (-1)^{g_s}Ze^{-\beta\frac{p^2}{2m}})|_0^\infty\right. \notag \\
    &+ \left.\frac{(-1)^{g_s}\beta Z}{3m}\int_0^\infty dp \, \frac{p^4e^{-\beta\frac{p^2}{2m}}}{1 + (-1)^{g_s}Ze^{-\beta\frac{p^2}{2m}}}\right) \tag{4.4.2} 
\end{align}
dove si è posto $p = \|\mathbf{p}\|$. Dopo qualche manipolazione algrebrica, possiamo riscrivere $(4.4.2)$ come segue
\begin{align}
    \frac{PV}{kT} &= \frac{2g_s\beta V}{3h^3}\int_0^\infty 4\pi p^2dp \, \frac{\frac{p^2}{2m}}{Z^{-1}e^{\beta\frac{p^2}{2m}} + (-1)^{g_s}} \notag \\
    &= \frac{2}{3}\beta U \tag{4.4.3}
\end{align}
ovvero 
\begin{align}
    PV = \frac{2}{3}U \tag{4.4.4}
\end{align}
La relazione $(4.4.4)$ resta valida anche se non è più valido il \textit{teorema di equipartizione dell'energia}. Questo perché ha fondamenti cinetici che sono più profondi.

\subsubsection*{Funzioni di Fermi}
Consideriamo ora il caso di fermioni e poniamo $x = \beta\frac{p^2}{2m}$, allora $p = \sqrt{\frac{2mx}{\beta}}$ e $dp = \frac{1}{2}\sqrt{\frac{2m}{\beta}}x^{-1/2}dx$. Sostituendo in $(4.4.1)$ otteniamo
\begin{align}
    \frac{PV}{kT} &= \frac{4\pi g_sV}{h^3}\int_0^\infty dx \, \frac{1}{2}\sqrt{\frac{2m}{\beta}}x^{-1/2}\frac{2m}{\beta}x\log(1 + Ze^{-x}) \notag \\
    &= \frac{4\pi g_sV}{h^3}\left(\frac{2m}{\beta}\right)^{3/2}\int_0^\infty dx \, x^{1/2}\log(1 + Ze^{-x}) \notag \\
    &= \frac{V}{\lambda^3}\frac{4}{\sqrt{\pi}}\int_0^\infty dx \, x^{1/2}\log(1 + Ze^{-x}) \tag{4.4.5}
\end{align}
dove 
\begin{align}
    \lambda = \left(\frac{\beta h^2}{2\pi m}\right)^{1/2} \notag
\end{align}
che contiene l'informazione sulla temperatura. Integrando per parti $(4.4.5)$ si ottiene
\begin{align}
    \frac{P}{kT} &= \frac{1}{\lambda^3}\frac{4}{\sqrt{\pi}}\left[\cancel{\left(\frac{2}{3}x^{3/2}\log(1 + Ze^{-x})\right)_0^\infty} + \frac{2Z}{3}\int_0^\infty dx \, \frac{x^{3/2}e^{-x}}{1 + Ze^{-x}} \right] \notag \\
    &= \frac{8}{3\lambda^3\sqrt{\pi}}\int_0^\infty dx \, \frac{x^{3/2}}{Z^{-1}e^x + 1} \tag{4.4.6}
\end{align}
Chiamiamo \textit{funzione di Fermi di ordine} $\frac{5}{2}$ il termine
\begin{align}
    f_{5/2}(Z) = \frac{8}{3\sqrt{\pi}}\int_0^\infty dx \, \frac{x^{3/2}}{Z^{-1}e^x + 1} \tag{4.4.7}
\end{align}
Più in generale, poniamo
\begin{align}
    f_\nu(Z) = \frac{1}{\Gamma(\nu)}\int_0^\infty dx \, \frac{x^{\nu - 1}}{Z^{-1}e^x + 1} \tag{4.4.8}
\end{align}
Per $Z << 1$, quindi al limite classico, si ha
\begin{align}
    f_{5/2}(Z) \sim Z \tag{4.4.9}
\end{align}
A questo punto possiamo riscrivere le equazioni caratteristiche dell'ensemble grancanonico quantistico termini delle funzioni di Fermi. In particolare
\begin{align}
    \mathcal{N} &= Z\frac{\partial}{\partial Z}\left(\frac{PV}{kT}\right)_{T, V} = \frac{V}{\lambda^3}Z\frac{\partial}{\partial Z}\left(f_{5/2}(Z)\right)_{T, V} \tag{4.4.10} \\
    U &= -\frac{\partial}{\partial \beta}\left(\frac{PV}{kT}\right)_{V, Z} = -Vf_{5/2}\frac{\partial}{\partial \beta}\left(\frac{1}{\lambda^3}\right)_{V, Z} \tag{4.4.11} 
\end{align}
Esplicitiamo i termini in $(4.4.10)$ e $(4.4.11)$. Per $\nu > 1$ si ha
\begin{align}
    Z\frac{\partial}{\partial Z}f_\nu(Z) &= \frac{Z}{\Gamma(\nu)}\int_0^\infty dx \, \frac{\partial}{\partial Z}\left(\frac{x^{\nu - 1}}{Z^{-1}e^x + 1}\right) \notag \\
    &= \frac{Z}{\Gamma(\nu)}\int_0^\infty dx \, \frac{x^{\nu - 1}}{(Z^{-1}e^x + 1)^2}\frac{e^x}{Z^2} \notag \\
    &= \frac{1}{\Gamma(\nu)}\int_0^\infty dx \, \frac{x^{\nu - 1}e^x}{Z(Z^{-1}e^x + 1)^2} \notag \\
    &= -\frac{1}{\Gamma(\nu)}\int_0^\infty dx \, x^{\nu - 1}\frac{d}{dx}\left(\frac{1}{Z^{-1}e^x + 1}\right) \notag \\
    &= -\frac{1}{\Gamma(\nu)}\left[\cancel{\left(\frac{x^{\nu - 1}}{Z^{-1}e^x + 1}\right)_0^\infty} - \int_0^\infty dx \, \frac{(\nu - 1)x^{\nu - 2}}{Z^{-1}e^x + 1}\right] \notag \\
    &= \frac{\nu - 1}{\Gamma(\nu)}\int_0^\infty dx \, \frac{x^{\nu - 2}}{Z^{-1}e^x + 1} \notag \\
    &= \frac{1}{\Gamma(\nu - 1)}\int_0^\infty dx \, \frac{x^{\nu - 2}}{Z^{-1}e^x + 1} \notag \\
    &= f_{\nu - 1}(Z) \tag{4.4.12}
\end{align}
Mentre
\begin{align}
    \frac{\partial}{\partial \beta}\left(\frac{1}{\lambda^3}\right) &= \frac{\partial}{\partial \beta}\left(\frac{\beta h^2}{2\pi m}\right)^{3/2} \notag \\
    &= \frac{3}{2}\left(\frac{h^2}{2\pi m}\right)^{3/2}\beta^{1/2} \notag
\end{align}
Pertanto, concludiamo che 
\begin{align}
    \frac{P}{kT} &= \frac{g_s}{\lambda^3}f_{5/2}(Z) \tag{4.4.13} \\
    \frac{\mathcal{N}}{V} &= \frac{g_s}{\lambda^3}f_{3/2}(Z) \tag{4.4.14} \\
    \frac{U}{V} &= \frac{3g_s}{2\beta}\frac{1}{\lambda^3}f_{5/2}(Z) \tag{4.4.15}
\end{align}
Inoltre, notiamo che è possibile riscrivere $f_\nu(Z)$ in termini di serie note, infatti
\begin{align}
    f_\nu(Z) &= \frac{1}{\Gamma(\nu)}\int_0^\infty dx \, \frac{x^{\nu - 1}}{Z^{-1}e^x + 1} \notag \\
    &= \frac{1}{\Gamma(\nu)}\int_0^\infty dx \, Zx^{\nu - 1}e^{-x}\sum_{l = 0}^\infty (-Ze^{-x})^l \notag \\
    &= \frac{1}{\Gamma(\nu)}\int_0^\infty dx \, x^{\nu - 1}\sum_{l = 0}^\infty (-1)^{l - 1}(Ze^{-x})^l \notag \\
    &= \sum_{l = 0}^\infty (-1)^{l - 1}\frac{Z^l}{l^\nu} \tag{4.4.16}
\end{align}
dove si è usato
\begin{align}
    \int_0^\infty dy \, y^{\nu - 1}e^{-ly} = \frac{\Gamma(\nu)}{l^\nu} \notag
\end{align}

\subsubsection*{Energia di Fermi}
Supponiamo ora di avere $\mathcal{N}$ particelle in un volume fissato $V$. Allora deve valere 
\begin{align}
    \frac{\mathcal{N}}{V} = \text{costante} \tag{4.4.17}
\end{align}
In particolare, confrontando $(4.4.17)$ con $(4.4.14)$, notiamo che, affinché la condizione $(4.4.17)$ sia soddisfatta, deve valere 
\begin{align}
    \frac{f_{3/2}(Z)}{\lambda^2} \propto 1 \tag{4.4.18}
\end{align}
dunque, dato che $\lambda^3$ va come $\beta^{3/2}$, otteniamo
\begin{align}
    f_{3/2}(Z) \propto \log Z^{3/2} \tag{4.4.19}
\end{align}
infatti $Z = e^{\beta\mu}$. Poniamo $\eta = \log Z$ (fissato) e supponiamo $\eta >> 1$, allora
\begin{align}
    f_{3/2}(Z) &= \frac{1}{\Gamma(3/2)}\int_0^\infty dx \, \frac{x^{1/2}}{Z^{-1}e^x + 1} \notag \\
    &= \frac{2}{\sqrt{\pi}}\int_0^\infty dx \, \frac{x^{1/2}}{e^{x - \eta} + 1} \notag \\
    &= \frac{2}{\sqrt{\pi}}\left[\cancel{\left(\frac{2}{3}\frac{x^{3/2}}{e^{x - \eta} + 1}\right)_0^\infty} + \frac{2}{3}\int_0^\infty dx \, \frac{x^{3/2}e^{x - \eta}}{(e^{x - \eta} + 1)^2}\right] \notag \\
    &= \frac{4}{3\sqrt{\pi}}\int_0^\infty dx \, \frac{x^{3/2}e^{x - \eta}}{(e^{x - \eta} + 1)^2} \tag{4.4.20}
\end{align}
L'integrale in $(4.4.20)$ è non nullo se e solo se $x \sim \eta$ (dato che $\eta >> 1$ per ipotesi), allora poniamo $t = x - \eta$ e sviluppiamo in serie di Taylor attorno a $\eta$ il termine $x^{3/2}$
\begin{align}
    f_{3/2}(Z) &\sim \frac{4}{3\sqrt{\pi}}\int_{-\eta}^\infty dx \, \frac{e^{x - \eta}}{(e^{x - \eta} + 1)^2}\left(\eta^{3/2} + \frac{3}{2}\eta^{1/2}(x - \eta) + \frac{3}{8}\eta^{-1/2}(x - \eta)^2 \dots\right) \notag \\
    &\sim \frac{4}{3\sqrt{\pi}}\int_{-\infty}^\infty dt \, \frac{e^t}{(e^t + 1)^2}\left(\eta^{3/2} + \frac{3}{2}\eta^{1/2}t + \frac{3}{8}\eta^{-1/2}t^2 \dots\right) \tag{4.4.21}
\end{align}
dove si è scambiato $-\eta$ con $-\infty$ nell'ipotesi $\eta >> 1$. Fissato $n \in \mathbf{N}$, poniamo
\begin{align}
    I_n = \int_{-\infty}^\infty dt \, t^n\frac{e^t}{(e^t + 1)^2} \tag{4.4.22}
\end{align}
Allora, sostituendo in $(4.4.21)$ si ha 
\begin{align}
    f_{3/2}(Z) \sim \frac{4}{3\sqrt{\pi}}\left(\eta^{3/2}I_0 + \frac{3}{2}\eta^{1/2}I_1 + \frac{3}{8}\eta^{-1/2}I_2 \dots\right) \tag{4.4.23}
\end{align}
Siamo ora interessati a determinare le proprietà del termine $I_n$ in base al valore dato di $n \in \mathbf{N}$- Notiamo che $I_n = 0$ se $n$ è dispari, infatti
\begin{align}
    (-t)^n\frac{e^{-t}}{(e^{-t} + 1)^2} &= -t^n\frac{e^{-t}}{e^{-2t} + 2e^{-t} + 1} \notag \\
    &= -t^n\frac{e^{-t}}{e^{-2t} + 2e^{-t} + 1}\left(\frac{e^{2t}}{e^{2t}}\right) \notag \\
    &= -t^n\frac{e^t}{1 + 2e^t + e^{2t}} \notag \\
    &= -t^n\frac{e^t}{(e^t + 1)^2} \notag
\end{align}
ovvero la funzione integranda è dispari. \\
Mentre, per $n$ pari è possibile ricondurre $(4.4.22)$ alla funzione $\zeta$ di Riemann. In particolare
\begin{align}
    I_n &= \int_{-\infty}^\infty dt \, t^n\frac{e^t}{(e^t + 1)^2} \notag \\
    &= 2\int_0^\infty dt \, t^n\frac{e^t}{(e^t + 1)^2} \notag \\
    &= -2\int_0^\infty dt \, \frac{\partial}{\partial \alpha}\left(\frac{t^{n - 1}}{e^{\alpha t} + 1}\right)_{\alpha = 1} \tag{4.4.24}
\end{align}
poniamo $u = \alpha t$, allora $t = \frac{1}{\alpha}u$ e $dt = \frac{1}{\alpha}du$. Sostituendo in $(4.4.24)$ si ottiene
\begin{align}
    I_n &= -2\frac{\partial}{\partial \alpha}\left(\frac{1}{\alpha^n}\int_0^\infty du \, \frac{u^{n - 1}}{e^u + 1}\right)_{\alpha = 1} \notag \\
    &= 2n\int_0^\infty du \, \frac{u^{n - 1}}{e^u + 1} \notag \\
    &= 2n\int_0^\infty du \, \frac{u^{n - 1}e^{-u}}{1 + e^{-u}} \notag \\
    &= 2n\int_0^\infty du \, u^{n - 1}e^{-u}\sum_{l = 0}^\infty (-e^{-u})^l \notag \\
    &= 2n\int_0^\infty du \, u^{n - 1}\sum_{l = 1}^\infty (-1)^{l - 1}(e^{-u})^l \notag \\
    &= 2n(n - 1)!\sum_{l = 1}^\infty \frac{(-1)^{l - 1}}{l^n} \tag{4.4.25}
\end{align}


\newpage
\subsection{Correzioni quantistiche alle equazioni caratteristiche}
Consideriamo un sistema di $N$ fermioni indipendenti e indistinguibili governato da un hamiltoniano di singola particella della forma
\begin{align}
    \mathcal{H}(\mathbf{p}) = c\|\mathbf{p}\| \tag{4.5.1}
\end{align}
Nel contesto dell'ensemble grancanonico quantistico si ha
\begin{align}
    U = \frac{4\pi g_sV}{h^3}\int_0^\infty dp p^2 \, \frac{pc}{Z^{-1}e^{\beta pc} + 1} \tag{4.5.2}
\end{align}
Poniamo $x = \beta pc$, allora $p = \frac{x}{\beta c}$ e $dp = \frac{1}{\beta c}dx$. Sostituendo in $(4.5.3)$ si ottiene 
\begin{align}
    U = \frac{4\pi g_sV}{\beta^4c^3h^3}\int_0^\infty dx \, \frac{x^3}{Z^{-1}e^x + 1} \tag{4.5.3} 
\end{align}
Facendo lo stesso per l'espressione del numero medio di particelle $\mathcal{N}$ si ottiene
\begin{align}
    \mathcal{N} = \frac{4\pi g_sV}{\beta^3c^3h^3}\int_0^\infty dx \, \frac{x^2}{Z^{-1}e^x + 1} \tag{4.5.4}
\end{align}
Sviluppiamo in serie $(4.5.4)$ per $Z << 1$
\begin{align}
    \mathcal{N} &= \frac{4\pi g_sV}{\beta^3c^3h^3}\int_0^\infty dx \, \frac{Zx^2e^{-x}}{1 + Ze^{-x}} \notag \\&\sim \frac{4\pi g_sV}{\beta^3c^3h^3}\left(Z\int_0^\infty dx \, x^2e^{-x} - Z^2\int_0^\infty dx \, x^2e^{-2x}\right) \notag \\
    &= \frac{4\pi g_s V}{\beta^3c^3h^3}(2Z - \frac{1}{4}Z^2) \tag{4.5.5}
\end{align}
dove si è usato 
\begin{align}
    \frac{1}{1 + Ze^{-x}} \sim 1 - Ze^{-x} \qquad Z << 1 \notag
\end{align}
Poniamo $n' = \frac{\mathcal{N}}{V}\frac{(\beta ch)^3}{4\pi g_s}$ ed esprimiamo $Z$ in funzione del numero medio di particelle. Per fare ciò è necessario risolvere la seguente equazione di secondo grado
\begin{align}
    \frac{1}{4}Z^2 - 2Z + n' = 0 \tag{4.5.6}
\end{align}
In particolare, $(4.5.6)$ ammette le seguenti soluzioni
\begin{align}
    Z_{\pm} = 4 \pm 2\sqrt{4 - n'} \sim 4 \pm \left(4 - \frac{1}{2}n' - \frac{1}{8}(n')^2\right) \tag{4.5.7}
\end{align}
dove si è usato $n' << 1$, dato che al limite classico $\frac{\mathcal{N}}{V} << 1$. La soluzione che ha senso fisicamente è quella minore al limite $n' \rightarrow 0$. In questo caso si tratta di $Z_-$, pertanto scegliamo
\begin{align}
    Z = \frac{1}{2}n' + \frac{1}{8}(n')^2 \tag{4.5.8}
\end{align}
Ora sviluppiamo anche $U$ in termini di $Z$
\begin{align}
    U &\sim \frac{4\pi g_sV}{\beta^4(ch)^3}\left(Z\int_0^\infty dx \, x^3e^{-x} - Z^2\int_0^\infty dx \, x^3e^{-2x}\right) \notag \\
    &= \frac{4\pi g_sV}{\beta^4(ch)^3}\left(6Z - \frac{3}{8}Z^2\right) \tag{4.5.9}
\end{align}
e sostituendo la soluzione ottenuta in $(4.5.8)$, tenendo conto solamente dei contributi fino all'ordine $2$, si ha
\begin{align}
    U &\sim \frac{4\pi g_sV}{\beta^4(ch)^3}\left(3n' - \frac{3}{4}(n')^2 - \frac{3}{16}(n')^2\right) \notag \\
    &= \frac{4\pi g_sV}{\beta^4(ch)^3}3n'\left(1 + \frac{5}{16}n'\right) \notag \\
    &= 3\mathcal{N}kT\left(1 + \frac{5}{16}n'\right) \tag{4.5.10}
\end{align}
Pertanto, dato che $PV = \frac{1}{3}U$, abbiamo ottenuto
\begin{align}
    PV = \mathcal{N}kT\left(1 + \frac{5}{16}n'\right) \tag{4.5.11}
\end{align}
ovvero la prima correzione alla pressione è positiva, ciò è dovuto al principio di esclusione di Pauli.









\newpage
\section*{Appendice}
\subsection*{Teorema di Liouville}
Consideriamo un sistema descritto da un'hamiltoniana $\mathcal{H} : \mathbf{R}^{6N} \longrightarrow \mathbf{R}$, dove $3N$ è il numero di gradi di libertà del sistema. L'hamiltoniana $\mathcal{H}$ definisce univocamente un campo di velocità $\mathbf{v} : \mathbf{R}^{6N} \longrightarrow \mathbf{R}^{6N}$ grazie alle equazioni del moto di Hamilton, ovvero
\begin{gather}
    \mathbf{v}(\mathbf{p}, \mathbf{q}) = \left(-\nabla_\mathbf{q}\mathcal{H}, \nabla_\mathbf{p}\mathcal{H}\right) \notag
\end{gather}
Si ha che $\nabla \cdot \mathbf{v} = 0$, infatti
\begin{gather}
    \nabla \cdot \mathbf{v} = \sum_{i = 1}^{3N}\left(-\frac{\partial^2 \mathcal{H}}{\partial p_i \partial q_i} + \frac{\partial^2 \mathcal{H}}{\partial q_i \partial p_i}\right) = 0 \notag
\end{gather}
per simmetria delle derivate seconde (l'hamiltoniana è sufficientemente regolare per scambiare l'ordine di derivazione per le derivate seconde). \\ \\
Sia $\Omega$ un (iper)volume contenuto nello spazio delle fasi, vogliamo dimostrare che vale la seguente identità
\begin{gather}
    \frac{d}{dt}m_{\Omega}(t) = \int_{\Omega} \nabla \cdot \mathbf{v} \, d\mathbf{p}d\mathbf{q} \notag
\end{gather}
ovvero che la variazione nel tempo del volume trasportato dal flusso di Hamilton è uguale al flusso di $\mathbf{v}$ entrante nel volume. Sia $\phi_t$ il flusso hamiltoniano, definito da
\begin{gather}
    \phi_t(\mathbf{p}, \mathbf{q}) = (\mathbf{p}(t), \mathbf{q}(t)) \notag
\end{gather}
$\phi_t$ associa ad ogni condizione iniziale $(\mathbf{p}, \mathbf{q})$ l'evoluzione temporale al tempo $t$ secondo le equazioni di Hamilton. In generale si ha
\begin{gather}
    m_\Omega(t) = \int_{\Omega(t)} d\mathbf{p}d\mathbf{q} \notag
\end{gather}
dove $\Omega(t) = \phi_t(\Omega)$. Applicando il cambio di variabili con il flusso si ottiene
\begin{gather}
    m_\Omega(t) = \int_\Omega \mathrm{Jac}(\phi_t(\mathbf{p}, \mathbf{q})) \, d\mathbf{p}d\mathbf{q} \notag
\end{gather}
Poniamo $A(\mathbf{p}, \mathbf{q}) = \mathrm{D}\phi_t(\mathbf{p}, \mathbf{q})$, dove
\begin{gather}
    \mathrm{D}\phi_t(\mathbf{p}, \mathbf{q}) = \left(\frac{\partial (\phi_t)_i}{\partial(\mathbf{p}, \mathbf{q})_j}\right)_{ij} \qquad i, j = 1, \dots, 6N \notag
\end{gather}
si ha che
\begin{gather}
    \frac{d}{dt}A = \mathrm{D}\mathbf{v}(\phi_t(\mathbf{p}, \mathbf{q})) \cdot \mathrm{D}\phi_t(\mathbf{p}, \mathbf{q}) = \mathrm{D}\mathbf{v}(\phi_t(\mathbf{p}, \mathbf{q}))A \notag
\end{gather}
dove si è usato
\begin{gather}
    \frac{d}{dt}\phi_t(\mathbf{p}, \mathbf{q}) = \frac{d}{dt}(\mathbf{p}(t), \mathbf{q}(t)) = \mathbf{v}(\phi_t(\mathbf{p}, \mathbf{q})) \notag
\end{gather}
ed è stato scambiato l'ordine delle derivate, supponendo sufficiente regolarità. Ora usiamo la seguente identità
\begin{gather}
    \frac{d}{dt}\det A = \det A \cdot \mathrm{tr}\left(A^{-1}\frac{d}{dt}A\right) \notag
\end{gather}
dunque, sostituendo $\mathrm{Jac}(\phi_t(\mathbf{p}), \mathbf{q}) = \det A$, si ottiene
\begin{align}
    \frac{d}{dt}m_\Omega(t) &= \int_\Omega \mathrm{Jac}(\phi_t(\mathbf{p}, \mathbf{q}))\cdot\mathrm{tr}(\mathrm{D}\mathbf{v}(\mathbf{p}, \mathbf{q})) \, d\mathbf{p}d\mathbf{q} \notag \\
    &= \int_{\Omega(t)} \nabla \cdot \mathbf{v} \, d\mathbf{p}d\mathbf{q} \notag
\end{align}
dove si è ritrasformato $\Omega$ in $\Omega(t)$ e quindi $\mathrm{Jac}(\phi_t(\mathbf{p}, \mathbf{q})) \, d\mathbf{p}d\mathbf{q} \rightarrow d\mathbf{p}d\mathbf{q}$ e si è esplicitata la traccia. In particolare se $\nabla \cdot \mathbf{v} = 0$ si ottiene
\begin{gather}
    \frac{d}{dt}m_\Omega(t) = 0 \notag
\end{gather}

\subsection*{Gamma di Eulero}
Definiamo la funzione Gamma di Eulero come segue
\begin{align}
    \Gamma(u) = \int_0^\infty e^{-t}t^{u - 1} \, dt \notag
\end{align}
definita per $u > 1$. Integrando per parti si ha
\begin{align}
    \Gamma(u) &= \int_0^\infty e^{-t}t^{u - 1} \, dt \notag \\
    &= \cancel{-e^{-t}t^{u - 1}|_0^\infty} + (u - 1)\int_0^\infty e^{-t}t^{u - 2} \, dt \notag \\
    &= (u - 1)\int_0^\infty e^{-t}t^{u - 2} \, dt \notag \\
    &= (u - 1)\Gamma(u - 1) \notag
\end{align}
osserviamo quindi che la funzione Gamma di Eulero è una generalizzazione del fattoriale per numeri reali maggiori di $1$. Per $u \in \mathbf{N}, u > 1$, iterando l'integrazione per parti si ottiene
\begin{align}
    \Gamma(u) = (u - 1)(u - 2)\dots(3)(2)\Gamma(1) \notag
\end{align}
dove
\begin{align}
    \Gamma(1) &= \int_0^\infty e^{-t} \, dt \notag \\
    &= -e^{-t}|_0^\infty = 1 \notag
\end{align}
pertanto $\Gamma(u) = (u - 1)!$. \\ \\
Vediamo ora come si comporta la Gamma di Eulero per $u$ abbastanza grande, riscriviamo $\Gamma(u)$ come segue
\begin{align}
    \Gamma(u + 1) &= \int_0^\infty e^{-t}e^{u\log t} \, dt \notag \\
    &= \int_0^\infty e^{-t + u\log t} \, dt \notag
\end{align}
Notiamo che la funzione integranda è scritta nella forma $e^{f(t)}$, dove
\begin{align}
    f(t) = -t + u\log t \notag
\end{align}
Determiniamo il punto di massimo per $f$, imponiamo $f'(t) = 0$, si ha
\begin{align}
    f'(t) &= -1 + \frac{u}{t} \notag
\end{align}
dunque $f'(t) = 0$ per $t = u$. Per la monotonia della funzione esponenziale si ha che il punto $t = u$ è anche punto di massimo per $e^{f(t)}$. Sviluppiamo $f$ in serie di Taylor attorno al suo punto di massimo
\begin{align}
    f(t) &\sim f(u) + \cancel{f'(u)(t - u)} + \frac{1}{2}f''(u)(t - u)^2 \notag \\
    &= -u + u\log u - \frac{1}{2u}(t - u)^2 \notag
\end{align}
notiamo che se usiamo lo sviluppo di $f$ attorno al suo punto di massimo si ha
\begin{align}
    e^{f(t)} &\sim e^{-u + u\log u - \frac{1}{2u}(t - u)^2} \notag \\
    &= e^{-u + u\log u}e^{-\frac{1}{2u}(t - u)^2} \notag
\end{align}
ovvero $e^{f(t)}$ diventa, a meno di un fattore costante, una gaussiana con varianza $\sigma = \sqrt{u}$. Dato che il punto di massimo della funzione integranda va come $u$, è valido, per $u$ abbastanza grande, usare lo sviluppo attorno al punto di massimo, dato che l'errore va come $u/\sqrt{u}$, che diventa infinitesimo per $u \rightarrow \infty$. Allora per $u$ abbastanza grande otteniamo
\begin{align}
    \Gamma(u + 1) &\sim e^{-u + u\log u}\int_0^\infty e^{-\frac{1}{2u}(t - u)^2} \, dt \notag \\
    &= e^{-u + u\log u}\int_{-u}^\infty e^{-\frac{s^2}{2u}} \, ds \notag \\
    &\sim e^{-u + u\log u}\int_{-\infty}^\infty e^{-\frac{s^2}{2u}} , ds \notag \\
    &= e^{-u + u\log u}\sqrt{2\pi u} \notag \\
    &= e^{-u}u^u\sqrt{2\pi u} \notag \\
    &\sim e^{-u}u^u \notag
\end{align}
dove si è posto $s = t - u$ e si è scambiamo l'estremo inferiore di integrazione $-u$ con $-\infty$, sempre nell'ipotesi $u$ abbastanza grande. Pertano otteniamo la formula di Stirling
\begin{align}
    \log\Gamma(u + 1) \sim u\log u - u \notag
\end{align}
per $u$ abbastanza grande.

\subsection*{Derivata di funzione integrale}
Consideriamo la seguente funzione integrale 
\begin{align}
    j(\alpha) = \int_{\alpha_0}^\alpha f(x, \alpha) \, dx \notag
\end{align}
dove $f$ è una funzione di classe $C^1$ nell'intervallo $[\alpha_0, \alpha]$. Per definizione, la derivata di $j$ è 
\begin{align}
    \frac{dj}{d\alpha} &= \lim_{h \rightarrow 0}\frac{1}{h}\left(\int_{\alpha_0}^{\alpha + h}f(x, \alpha + h) \, dx - \int_{\alpha_0}^\alpha f(x, \alpha) \, dx\right) \notag
\end{align}
Aggiungendo e sottraendo un termine
\begin{align}
    \frac{1}{h}\int_{\alpha_0}^{\alpha + h}f(x, \alpha) \, dx \notag
\end{align}
si ottiene 
\begin{align}
    \frac{dj}{d\alpha} &= \lim_{h \rightarrow 0}\frac{1}{h}\left(\int_{\alpha_0}^{\alpha + h}(f(x, \alpha + h) - f(x, \alpha)) \, dx + \int_{\alpha}^{\alpha + h}f(x, \alpha) \, dx\right) \notag \\
    &= \int_{\alpha_0}^\alpha \frac{\partial f}{\partial \alpha} \, dx + f(\alpha, \alpha) \notag
\end{align}
dove si è applicato il \textit{teorema della media integrale} per entrambi i termini. 

\subsection*{Limite gaussiano della distribuzione di Poisson}
La distribuzione di Poisson è data da 
\begin{align}
    P(x) = e^{-\lambda}\frac{\lambda^x}{x!} \notag
\end{align}
per qualche opportuno $\lambda \in \mathbf{R}$ e per ogni $x$ nell'insieme di definizione. Supponiamo 
\begin{align}
    \Delta = x - \lambda >> 1 \qquad \frac{\Delta}{\lambda} << 1 \notag
\end{align}
Per la formula di Stirling possiamo scrivere 
\begin{align}
    x! \sim \sqrt{2\pi}e^{-x}x^{x + 1/2} \notag
\end{align}
pertanto 
\begin{align}
    P(x) \sim \frac{1}{\sqrt{2\pi}}e^{-\lambda}\frac{\lambda^x}{e^{-x}x^{x + 1/2}} = \frac{1}{\sqrt{2\pi}}e^{-\lambda + x\log \lambda + x - (x + 1/2)\log x} \notag
\end{align}
Consideriamo solamente l'esponente, sostituendo $x = \lambda + \Delta$ si ha 
\begin{align}
    \cancel{-\lambda} + (\lambda + \Delta)\log \lambda + (\cancel{\lambda} + \Delta) + (\lambda + \Delta + 1/2)\log(\lambda + \Delta) \notag
\end{align}


\end{document}
